\documentclass{article}

% NeurIPS 2026 preprint form: author names visible, no line numbers,
% no "Submitted to ..." footer. For the anonymized double-blind
% submission package use \usepackage{neurips_2026} without [preprint];
% for the camera-ready, use \usepackage[final]{neurips_2026}.
\usepackage[preprint]{neurips_2026}

\usepackage[utf8]{inputenc}
\usepackage[T1]{fontenc}
\usepackage{newunicodechar}
\usepackage{hyperref}
\usepackage{url}
\usepackage{booktabs}
\usepackage{longtable}     % pandoc renders the survey/contrast tables as longtables
\usepackage{array}
\usepackage{calc}          % provides \real{} used in pandoc longtable column widths
\usepackage{amsfonts}
\usepackage{amssymb}
\usepackage{amsmath}
\usepackage{enumitem}
\usepackage{nicefrac}
\usepackage{microtype}
\usepackage{xcolor}
\usepackage{graphicx}

% --- pandoc-fragment support macros (the body is a pandoc fragment, so the
% helper definitions pandoc normally puts in its standalone template must be
% supplied here) ---
% Caption-less longtables: pandoc emits \def\LTcaptype{none}; define the counter
% so longtable's \refstepcounter{\LTcaptype} is a harmless no-op under pdflatex.
\makeatletter
\@ifundefined{c@none}{\newcounter{none}}{}
\makeatother
% Tight (compact) list spacing used by pandoc.
\providecommand{\tightlist}{\setlength{\itemsep}{0pt}\setlength{\parskip}{0pt}}

% This NeurIPS version is a format-only port of the arXiv "short" version
% (../../arxiv/frontier_and_localhost/short/). The body is generated from that
% Markdown via pandoc --natbib (\citep/\citet), so we preserve the source's
% manual section numbering rather than re-numbering: LaTeX adds no section
% numbers of its own, and the numbers shown ("3.5", "Appendix G.8", ...) come
% from the heading text. This keeps every literal "S3.5 / Appendix G" cross
% reference in the prose valid with zero renumbering.
\setcounter{secnumdepth}{0}

% Non-ASCII glyphs the pandoc body uses, mapped for pdflatex.
\newunicodechar{§}{\S{}}
\newunicodechar{×}{\ensuremath{\times}}
\newunicodechar{→}{\ensuremath{\rightarrow}}
\newunicodechar{⊂}{\ensuremath{\subset}}
\newunicodechar{†}{\textdagger{}}
\newunicodechar{≥}{\ensuremath{\geq}}
\newunicodechar{−}{\ensuremath{-}}
\newunicodechar{∘}{\ensuremath{\circ}}

\title{Frontier and Localhost:\\How Production AI Learns Outside the Weights}

\author{%
  Mikhail L. Arbuzov\\
  Independent Researcher\\
  \texttt{mike.arbuzov54@gmail.com}\\
  \And
  Lee Mosbacker\\
  Independent Researcher\\
  \texttt{lee.mosbacker@gmail.com}\\
  \And
  Sisong Bei\\
  Independent Researcher\\
  \texttt{qurining@gmail.com}\\
  \AND
  Ziwei Dong\\
  Independent Researcher\\
  \texttt{ziwei.dong@alumni.emory.edu}\\
  \And
  Dmitri Kalaev\\
  Independent Researcher\\
  \texttt{kalaevdr@gmail.com}\\
  \And
  Alexey A. Shvets\\
  Palo Alto Networks\\
  \texttt{ashvets@paloaltonetworks.com}\\
}

\begin{document}

\maketitle

\begin{abstract}
Production LLM systems increasingly adapt outside the model weights. After deployment failures, teams modify prompts, rules, memories, skills, tools, eval suites, routing graphs, and governance pipelines on a cadence that frontier weight updates cannot match. But this scaffold layer is still mostly maintained as \emph{patchwork}: fixes are proposed by intuition, committed with weak credit assignment, accumulated without pruning, and promoted beyond the scope where they were validated. This paper formalises a disciplined alternative, which we call \emph{artifact-layer descent}. The pipeline is concrete. Recurring residuals are assigned to scaffold coordinates. Candidate artifact deltas are tested against patch loss. Accepted deltas persist with rollback. Promotion across contexts is bounded by \emph{evidence radius}: a local fix spreads only as far as the evidence supports. Surveying approximately 130 production and research systems from 2023 to 2026, we find current systems partially instantiate this loop but leave a central architecture gap, namely governed scaffold optimization across contexts. The contribution is not the claim that prompts are weights. It is to name the missing optimizer for an adaptation layer the field has already built.
\end{abstract}

\subsection{1. The gradient moved outside the
model}\label{the-gradient-moved-outside-the-model}

The standard mental model of an LLM system is the model. Pre-training
fixes the weights, fine-tuning nudges them, deployment wraps a thin
prompt-plus-tool layer around the result and ships. Production in 2025
and 2026 has stopped behaving this way. Frontier models update on cycles
measured in months; the systems built on top of them update
continuously, in markdown files, persistent memories, retrieved skills,
registered tools, evaluation suites, agent topologies,
version-controlled prompt repositories, PR-gated governance. A growing
share of deployment-time adaptation now occurs outside the weights, in a
fast-changing \emph{local scaffold} wrapped around a slow-changing
frontier. The section title is a metaphor. The gradient inside the model
still happens during training, and the scaffold layer is not a gradient
in the calculus sense. The substantive claim is the milder one.
Production AI has accidentally created an external learning surface, and
right now it is being operated as \emph{patchwork} rather than as a
disciplined optimization process.

Read together, these artifacts form a single external adaptation surface
(the six substrates catalogued in §4) that current production practice
still treats as scattered engineering. A Claude Skill is a learned
procedural weight kept outside the model. A Cursor Rule is a local
parameter for one project's distribution shift. A RAG connector is
patch-specific memory access. A tool registry is capability provisioning
along an axis the base model cannot reach. A PR review against a shared
rule repository is a promotion gate on a federated update. An eval suite
that blocks merge is a validation-loss check before promotion. A
cross-tenant artifact repository is \emph{federated artifact promotion}
across distributed shards: scope expansion under a gate, not arithmetic
averaging. None of these mechanisms is new; what is new is naming the
surface and asking what optimization discipline it needs --- and
visibly lacks.

This paper surveys approximately 130 production and research systems,
identifies the six substrates where the scaffold lives, formalises the
two-loop architecture (local update, cross-context promotion under a
gate), and names a missing architecture as the most consequential open
problem: governed cross-tenant scaffold optimization with versioned
lineage, role-based access, and rollback. The single-line thesis is
this: \textbf{the field has built an external adaptation layer for LLMs,
but not the optimizer that should govern it.} The rest of the paper
develops that observation into a survey of where production sits on a
\emph{patchwork to random-search to descent} gradient, a formalisation
of the disciplined limit case under stated assumptions, and a named
missing architecture where the optimizer most plainly does not yet
exist.

\subsection{2. Why patch errors live outside the
weights}\label{why-patch-errors-live-outside-the-weights}

The scaffold is not merely where production systems happen to store
patches. For high-churn, tenant-specific, auditable adaptation, the
scaffold is \emph{often the practical substrate}. The reason is not that
no alternative exists. It is that under deployment-time constraints
(release cadence, tenant isolation, inspectability, rollback), the
scaffold is cheap, locally scoped, reversible, and updatable at human
time-scales while the weights are not.

\textbf{Frontier training optimises for breadth.} The weights are tuned
for cross-patch generalisation. They smooth over local conventions,
contradictory tenant preferences, and rare workflow-specific failures in
order to preserve broad competence across millions of unseen
deployments. The smoothing is the training objective, not an artefact.
RLHF reduces output diversity and homogenises preferred styles across
users \citep{kirk2024understanding, padmakumar2023diversity}, and
instruction-tuning data composition rewards patterns that generalise
across the annotator pool rather than patterns specific to any one
deployment. The frontier model is therefore \emph{trained to be unable}
to encode a single tenant's idiosyncrasies as preferred behaviour.
\citet{tiwari2026fastandslow} supply an independent two-timescale
optimisation motivation from the in-context / in-weights side:
in-context adaptation is cheap and rapid but capacity-bounded,
in-weights adaptation is expressive but induces catastrophic forgetting
and plasticity loss; the scaffold layer is the substrate that absorbs
what in-context adaptation can carry cheaply, leaving the weights free
to generalise.

\textbf{Production reliability is achieved in depth.} Inside a single
deployment patch, reliability is determined by exactly what global
training had to smooth away: local naming conventions, team-specific
workflows, idiosyncratic tool chains, customer-specific risk tolerances,
recurring local mistakes, hidden evaluator expectations. The per-patch
quantities from our previous work are local by construction, and a
library calibrated to one patch under-covers the next.

\textbf{Encoding every patch in the weights is structurally infeasible.}
Three constraints rule out the obvious approach. \emph{Economic:}
millions of patches multiplied by billions of parameters times per-patch
fine-tuning cost is not a viable training-compute budget.
\emph{Release-cycle:} the patch wants daily updates; weight releases
ship quarterly at best. \emph{Cross-patch interference:} overfitting the
weights to one tenant's preferences degrades adjacent tenants whose
preferences contradict. The result is brittle in exactly the way the
scaffold is not.

Per-tenant fine-tunes, LoRA adapters, distilled small models, and
learned routing each absorb some patch residuals in principle. In
practice the scaffold dominates, because it is the substrate that
satisfies all four production constraints at once: cheap, inspectable,
reversible, locally scoped, and updatable on a daily cadence rather than
a release cadence. The placement question that \emph{Architecture of
Errors} left open (where do the interventions live?) has this as its
answer. The frontier model generalises. The local scaffold specialises.
Reliability comes from governing that specialisation.

\textbf{Three stages of scaffold maintenance recur across the corpus.}
They differ in how candidate updates are produced, whether credit
assignment is explicit, and whether promotion is evidence-bounded. Most
production deployments sit at Stage 1 or 2. The math of §3 applies to
Stage 3.

{\def\LTcaptype{none} % do not increment counter
\begin{longtable}[]{@{}
  >{\raggedright\arraybackslash}p{(\linewidth - 8\tabcolsep) * \real{0.2000}}
  >{\raggedright\arraybackslash}p{(\linewidth - 8\tabcolsep) * \real{0.2000}}
  >{\raggedright\arraybackslash}p{(\linewidth - 8\tabcolsep) * \real{0.2000}}
  >{\raggedright\arraybackslash}p{(\linewidth - 8\tabcolsep) * \real{0.2000}}
  >{\raggedright\arraybackslash}p{(\linewidth - 8\tabcolsep) * \real{0.2000}}@{}}
\toprule\noalign{}
\begin{minipage}[b]{\linewidth}\raggedright
Stage
\end{minipage} & \begin{minipage}[b]{\linewidth}\raggedright
Update mechanism
\end{minipage} & \begin{minipage}[b]{\linewidth}\raggedright
Credit assignment
\end{minipage} & \begin{minipage}[b]{\linewidth}\raggedright
Promotion discipline
\end{minipage} & \begin{minipage}[b]{\linewidth}\raggedright
Dominant failure mode
\end{minipage} \\
\midrule\noalign{}
\endhead
\bottomrule\noalign{}
\endlastfoot
\textbf{1. Patchwork} & Intuition-driven local edits & Implicit / ad hoc
& None & Bloat \\
\textbf{2. Random artifact search} & Many variants tried under eval &
Weak; eval is the only signal & Eval-gated CI but no
failure-to-coordinate tracing & Overfitting to in-distribution patch \\
\textbf{3. Artifact-layer descent} & Residual to coordinate to delta &
Explicit per failure mode & Evidence-radius-bounded; reversible & (the
discipline this paper formalises) \\
\end{longtable}
}

The paper's central claim is not that today's production systems already
run this loop. Most do not. The claim is that a patchwork loop
\emph{becomes} artifact-layer descent once the five conditions of §3 are
in place. The survey of §§5--8 asks which deployed systems have closed
which of them; §9 names the one closure no public system has yet
reached.

\subsection{3. Artifact-layer descent}\label{artifact-layer-descent}

\textbf{The disciplined version of scaffold maintenance is one loop, run
again and again over the life of a deployment.} It does not react to
single failures; it works from their accumulated record, and it ends
with a decision about how widely to trust each fix. As the deployment
runs, it logs what went wrong --- failed outputs, user corrections,
rejected answers, eval regressions, tool-call errors. Reflecting over
that record, the loop clusters the failures into recurring \emph{failure
modes} and takes up the ones that matter most: a mode earns attention in
proportion to how often it recurs and how badly it fails ---
repeatability times severity, its contribution to the deployment's
failure rate. For such a mode, locate the single scaffold
coordinate responsible --- a stale memory entry, an under-specified
instruction, a missing tool argument, a brittle retrieval rule, a
misrouted agent edge --- and synthesise one candidate edit to it.
Measure whether the edit lowers \emph{patch loss} --- the deployment's
failure rate on its own recurring tasks --- and let a gate accept or
reject it on that measurement. Keep an accepted edit in the scaffold;
prune it or roll it back if a later check shows it made things worse.
Then take the step most systems skip: decide how far the edit should
travel, and promote it to other deployments only as far as it has
actually been validated. That last limit --- a fix earned on one
deployment spreads only as far as its evidence reaches --- is the edit's
\emph{evidence radius}. The rest of this section makes each step
precise; §§4--9 then ask which deployed systems have closed which of
them.

\textbf{Most current production does not close this loop.} A failure
happens, someone adds a prompt or rule or memory entry, the artifact
pile grows, nobody knows what caused improvement, nobody knows when to
prune, and fixes leak into broader contexts. This is patchwork (Stage 1
of §2), not stochastic descent. Even systems with eval-gated CI (Stage
2) usually lack failure-to-coordinate credit assignment, so improvements
are real but unattributed and brittle to transfer. Scaffold updates
\emph{become} stochastic descent (Stage 3) only when five pieces are in
place: a measurable patch loss; captured residuals; a failure traced
back to the scaffold coordinate that caused it; a synthesised candidate
edit; and a gate that checks whether the edit lowers loss before it is
kept or promoted. The descent picture below is mechanism when those five
conditions hold, and modeling language otherwise.

\textbf{Two loops, in plain English first.} \emph{Loop 1} is the local
fix: a recurring failure type produces a candidate edit, the edit is
tested, and the accepted edit changes one deployment's scaffold.
\emph{Loop 2} is the separate decision of whether that local fix should
be shared with other deployments --- promoting it from one repository to
a team-wide template, from one tenant's playbook to a cross-tenant skill
library, and so on. The two loops turn different knobs. Loop 1 improves
the scaffold against a single deployment's own failures; Loop 2 changes
\emph{which} deployment's failures the scaffold is being tuned against.

\textbf{Assumptions and scope of the formalization.} The descent picture
is only as good as the conditions it rests on, so we state them before
formalising it; where a condition fails, the same picture should be read
as analogy rather than mechanism. The argument below holds when: (i) the
frontier model's weights stay frozen while the scaffold is being
adapted; (ii) each scaffold edit is a persistent, inspectable artifact,
not a transient prompt that vanishes after one session; (iii) the
deployment exposes a loss you can actually measure --- a failure rate, a
correction rate, an eval score; (iv) the process that proposes edits
can, at least sometimes, trace a failure back to the scaffold coordinate
responsible for it; (v) the gate is a noisy estimator of whether an edit
helps, not an oracle that knows the true change in loss; and (vi) Loop-2
promotion widens an artifact's scope under a gate, rather than
numerically averaging comparable updates the way weight-space federated
learning does. Survey claims of the form ``no surveyed system has X''
are bounded by the audited public corpus and the survey cutoff of May
2026.

\textbf{Patch scope hierarchy.} \emph{Patch} is used throughout this
paper with a defined scope hierarchy: USER ⊂ PROJECT ⊂ STACK ⊂ TENANT ⊂
CORE. Throughout, \emph{patch} denotes a deployment scope, never an
individual fix; the undisciplined practice of editing the scaffold by
hand is \emph{patchwork}. A USER patch is one human's account or fork. A PROJECT is one
repository, workflow, or product workspace. A STACK is a curated bundle
of projects under one team. A TENANT is an organisation or customer.
CORE is the cross-tenant universal scope. Each level has its own loss
--- the average failure rate over all the patches it contains. When the
text says ``cross-tenant promotion'' it specifically means moving
artifacts up to the TENANT or CORE level. ``Patch-specific adaptation''
without qualifier means descent on a USER or PROJECT-level loss.

The patch loss is the expected residual failure rate of the
model-scaffold composite on patch \(D\):

\[L_D(\theta_s) = \mathbb{E}_{x \sim D}\!\left[\ell(\pi(\theta_M, \theta_s, x))\right],\]

where \(\theta_s\) is a structured scaffold parameter object over six
substrate coordinates and \(\ell\) is an operational loss (failure
indicator, user-correction rate, eval error, schema violation, reviewer
disagreement). The important distinction is substrate. \(\theta_M\) is
fixed during deployment. \(\theta_s\) is the object being fit.

A failure does not yield an analytic gradient. It yields a candidate
edit --- synthesised by reflection over the failure history, and
proposed in order of the recurrence × severity of the mode it addresses.
The gate \(G_t\) accepts the edit if its noisy estimate of the
directional loss change clears a margin. The accepted edit applies to
the current scaffold; rejected edits leave the scaffold unchanged:

\[\theta_s^{t+1} = \begin{cases} \theta_s^t \oplus z_t, & G_t(z_t)=1, \\ \theta_s^t, & G_t(z_t)=0, \end{cases}\]

with the gate accept condition

\[G_t(z_t)=1 \iff \widehat{\Delta}_{z_t} L_D(\theta_s^t) + \lambda\,\Omega(z_t) \leq -\tau_t,\]

where \(\widehat{\Delta}\) is the gate's estimator of the loss change,
\(\Omega \geq 0\) is a complexity penalty on the edit (instruction
length, tool-permission breadth, promotion scope), \(\lambda \geq 0\)
regularises that penalty, and \(\tau_t \geq 0\) is the required
improvement margin. Production systems implement this without writing
the equation: an eval blocks a merge on regression, a reviewer rejects a
vague rule, a governance system prevents a local artifact from being
shared until evidence accumulates. Different surveyed gate families
instantiate the abstraction differently. Eval-cascade, RL-threshold,
LLM-judge, surrogate verifier, and reviewer-judgment families directly
populate the per-edit gate. Semantic-merge and MAP-Elites archive
selection act as proposal-pruning steps that \emph{precede} the per-edit
gate. The per-system mapping is in Appendix C.

\textbf{Two loops, two distinct knobs.} Consider a concrete case. An
engineer notices that a coding agent keeps misnaming a class in one
repository. The fix is a one-line rule in that repo's instruction file.
The edit is small. Its scope is one project. The same engineer could
promote the rule to a team-wide template, or to the company's CORE
instructions, or to a cross-tenant skill registry. The \emph{edit} did
not change. The \emph{scope} did. And the gate that should let it
through gets stricter at each step, because the loss it is being tested
against is now an average over more patches, and a fix that helps one
repo can hurt another.

That observation maps onto two distinct knobs in the optimisation
analogy. The
first is \emph{edit magnitude} --- how much a single accepted edit
changes the local scaffold. This is the analogue of a learning rate: it
sets how big each step is. The second is \emph{promotion radius} --- how
many downstream patches the accepted edit is applied to. This does not
change the step size; it changes the \emph{target} the step is aimed at,
that is, which average failure rate the scaffold is being tuned to
reduce. Accept an edit at the USER level and you are reducing the
failure rate on one patch (\(L_D\)). Promote it to a PROJECT and you are
now reducing the average failure rate across all of that project's
patches; promote it to TENANT or CORE and the target becomes a federated
average over many tenants' patches. A larger radius therefore aims the
same edit at a broader, harder-to-please target, and an edit that lowers
the narrow average can raise the broad one. The two knobs are set
independently, and the gate must be calibrated against both. This is why
governance is not administrative overhead. It is the regularization
system that keeps a small fix from being aimed at the wrong target.

Loop 2 is \emph{federated artifact promotion} --- selection, merge,
curation, review, distribution --- not arithmetic averaging of numeric
updates; §9 sharpens the FedAvg analogy and its limits.

Each loop has a gate. Gates are \textbf{Auto-gated} when the decision is
an operationalised algorithmic criterion (RL threshold, eval score,
LLM-judge with quantified agreement, surrogate verifier, benchmark
cascade). They are \textbf{Human-gated} when the decision is reviewer
judgment. They are \textbf{None} when the loop is absent. These five
terms (Loop 1, Loop 2, Auto-gated, Human-gated, None) anchor §§5--9.

The descent argument is \emph{mechanism} when these five conditions hold
and \emph{metaphor} when they are absent. Appendix E formalises the
proof assumptions A1--A5 --- distinct from the operating assumptions
(i)--(vi) above --- needed for a one-step descent inequality and a
convergence proposition of the loss to a \(G\)-stable scaffold, and
works through the RAG, pitch-review, tool-use, and memory cases that
make the mechanism visible.

\subsection{3.5. SkillOpt and the convergence on artifact-layer
descent}\label{skillopt-and-the-convergence-on-artifact-layer-descent}

Recent work has begun to operationalise all five §3 conditions for
skill-document adaptation, and SkillOpt \citep{yang2026skillopt} is the
clearest case: it treats a single markdown skill document as the
trainable external state of a frozen agent and edits it with seven
explicit deep-learning-style controls --- rollout batches; reflection
minibatches; add/delete/replace edits; a textual learning-rate budget
with cosine decay; a held-out selection-split accept gate; a
rejected-edit buffer that keeps failed proposals as negative feedback;
and an epoch-wise slow/meta update for longer-horizon consolidation. Its
framing --- \emph{skill optimization as deep-learning-style optimization
over an external natural-language state} --- is exactly the disciplined
gradient-like update mechanism §3 argues is missing in current
production. SkillOpt is not an isolated proof of concept but the
cleanest member of a rapidly crystallising space: the same convergence
pressure this paper identifies is arriving from at least 14 independent
directions in 2025--2026, each instantiating a different subset of §3's
five conditions and targeting a different substrate (skill document,
harness code, system prompt, shared memory).

Mapped onto §3, SkillOpt instantiates each abstract quantity with a
concrete mechanism; Appendix G expands every row with the exact
algorithm-level mechanism, including pointers to SkillOpt's Algorithm 1
line numbers.

{\def\LTcaptype{none} % do not increment counter
\begin{longtable}[]{@{}
  >{\raggedright\arraybackslash}p{(\linewidth - 6\tabcolsep) * \real{0.2500}}
  >{\raggedright\arraybackslash}p{(\linewidth - 6\tabcolsep) * \real{0.2500}}
  >{\raggedright\arraybackslash}p{(\linewidth - 6\tabcolsep) * \real{0.2500}}
  >{\raggedright\arraybackslash}p{(\linewidth - 6\tabcolsep) * \real{0.2500}}@{}}
\toprule\noalign{}
\begin{minipage}[b]{\linewidth}\raggedright
§3 quantity
\end{minipage} & \begin{minipage}[b]{\linewidth}\raggedright
SkillOpt instantiation
\end{minipage} & \begin{minipage}[b]{\linewidth}\raggedright
Mechanism
\end{minipage} & \begin{minipage}[b]{\linewidth}\raggedright
Numerical signature
\end{minipage} \\
\midrule\noalign{}
\endhead
\bottomrule\noalign{}
\endlastfoot
\(\theta_s\) (scaffold) & \texttt{best\_skill.md} & persistent text
artifact, exported as the deployment object & 379--1,995 tokens;
×2.5--53 growth \\
\(L_D\) (patch loss) & held-out hard / exact-match score & per-benchmark
evaluator & 6 benchmarks \\
\(Q_t\) (proposal) & optimizer-model reflection over failure and success
minibatches & LLM-judge over rollout batch & \(B_m{=}8\), parallel
reflectors \\
\(z_t\) (edit) & add/delete/replace patch on the skill & structured JSON
edit & 1--4 accepted edits per skill \\
\(L_t\) (textual learning rate) & bounded edit budget per step &
constant/linear/cosine/autonomous schedule &
\(L_t \in \{1,2,4,8,16\}\) \\
\(G_t\) (gate) & held-out selection-split accept gate & strictly-greater
validation & ties rejected \\
rejected-edit buffer & epoch-local store of failed proposals &
training-time negative feedback only & zero inference cost \\
slow/meta update & epoch-wise longitudinal guidance written to a
protected slow-update field & second-order consolidation & \(-22.5\) pts
on SpreadsheetBench when removed \\
\end{longtable}
}

The headline empirics make the descent picture testable rather than
aspirational. Across six benchmarks, seven target models, and three
execution harnesses (direct chat, Codex, Claude Code), SkillOpt is best
or tied-best on all 52 evaluated (model, benchmark, harness) cells,
lifting GPT--5.5 by +23.5 / +24.8 / +19.1 points over no-skill under
direct chat / Codex / Claude Code, and beating the strongest per-cell
baseline drawn from TextGrad, GEPA, Trace2Skill, EvoSkill, human-skill,
and one-shot-LLM by +5.4 points on average. Transfer is also measured: a
SpreadsheetBench skill trained inside the Codex loop transfers to Claude
Code at +59.7 over that harness's no-skill baseline, and an
OlympiadBench skill yields positive gains on Omni-MATH at every model
scale tested. In §3 vocabulary: SkillOpt closes Loop 1 under an
algorithmic gate, runs an epoch-wise auto-gated stabiliser as a second
algorithmic loop, and reports cross-model / cross-harness /
cross-benchmark transfers that read as \emph{empirical evidence-radius
measurements} --- Loop-2 scope expansion that is sound for the
demonstrated radii while no cross-organisational governance is shipped.

\textbf{The convergent landscape.} Five clusters of concurrent work
instantiate complementary subsets of §3's mechanism stack, each
targeting a different substrate or exposing a different constraint:
\emph{reflective-evolution}
\citetext{\citealp[GEPA,][]{agrawal2025gepa}; \citealp{tiwari2026fastandslow}},
\emph{reflective context} \citep{vassilyev2026rcl}, \emph{harness
optimisation}
\citetext{\citealp[Meta-Harness,][]{lee2026metaharness}; \citealp[AHE,][]{lin2026ahe}; \citealp[MOSS,][]{cai2026moss}; \citealp[CANTANTE,][]{zehle2026cantante}; \citealp{ong2026harnesseval}},
\emph{multi-objective + lifecycle}
\citetext{\citealp[MOCHA,][]{tanjim2026mocha}; \citealp[Ratchet,][]{zhang2026ratchet}},
and \emph{cross-task transfer}
\citetext{\citealp[ICT,][]{shalev2026ict}; \citealp[Combee,][]{li2026combee}; \citealp[CORAL,][]{qu2026coral}; \citealp[PACE,][]{ling2026pace}}.
The table below indexes what each instantiates in §3 and what it leaves
open; Appendix G walks each cluster in detail with its §3 mapping stated
explicitly.

{\def\LTcaptype{none} % do not increment counter
\begin{longtable}[]{@{}
  >{\raggedright\arraybackslash}p{(\linewidth - 6\tabcolsep) * \real{0.2500}}
  >{\raggedright\arraybackslash}p{(\linewidth - 6\tabcolsep) * \real{0.2500}}
  >{\raggedright\arraybackslash}p{(\linewidth - 6\tabcolsep) * \real{0.2500}}
  >{\raggedright\arraybackslash}p{(\linewidth - 6\tabcolsep) * \real{0.2500}}@{}}
\toprule\noalign{}
\begin{minipage}[b]{\linewidth}\raggedright
Cluster
\end{minipage} & \begin{minipage}[b]{\linewidth}\raggedright
Representative systems
\end{minipage} & \begin{minipage}[b]{\linewidth}\raggedright
§3 mechanism instantiated
\end{minipage} & \begin{minipage}[b]{\linewidth}\raggedright
What is left open
\end{minipage} \\
\midrule\noalign{}
\endhead
\bottomrule\noalign{}
\endlastfoot
Reflective-evolution & GEPA; Tiwari et al. & \(Q_t\) over trajectory
evidence; two-timescale motivation & persistent \(\theta_s\); \(L_t\)
budget; held-out \(G_t\) \\
Reflective context & RCL (Vassilyev et al.) & \(\theta_s\) over context
space; \(\Omega(z)\) as generalisation penalty & cross-deployment Loop
2 \\
Harness optimisation & Meta-Harness; AHE; MOSS; CANTANTE; Ong et al. &
credit assignment over S5; source-level extension & governed Loop 2
across orgs \\
Multi-objective + lifecycle & MOCHA; Ratchet & constraint Pareto
frontier; hygiene-as-precondition & gradient-analogue closure \\
Cross-task transfer & ICT; Combee; CORAL; PACE & empirical
evidence-radius; async multi-agent Loop 2 within one org &
cross-organisational gate, lineage, rollback \\
\end{longtable}
}

\textbf{What the convergence establishes.} Read against §3's mechanism
stack and §9's four-property criterion, the family establishes four
things: the scaffold-versus-weights decomposition is independently
confirmed; credit assignment is the central engineering bottleneck;
lifecycle hygiene is a precondition for descent rather than an
alternative to it; and no member combines Loop-1 automation with
cross-organisational Loop-2 governance --- so §9's gap is confirmed, not
reduced, by the convergence. Appendix G develops each claim, gives the
full SkillOpt correspondence and its contrast against the six measured
baselines, and walks each convergent cluster.

\subsection{4. Six scaffold substrates}\label{six-scaffold-substrates}

A patch-local scaffold composes six substrates. They are the coordinates
of \(\theta_s\), each independently editable.

{\def\LTcaptype{none} % do not increment counter
\begin{longtable}[]{@{}
  >{\raggedright\arraybackslash}p{(\linewidth - 4\tabcolsep) * \real{0.3333}}
  >{\raggedright\arraybackslash}p{(\linewidth - 4\tabcolsep) * \real{0.3333}}
  >{\raggedright\arraybackslash}p{(\linewidth - 4\tabcolsep) * \real{0.3333}}@{}}
\toprule\noalign{}
\begin{minipage}[b]{\linewidth}\raggedright
Substrate
\end{minipage} & \begin{minipage}[b]{\linewidth}\raggedright
What persists
\end{minipage} & \begin{minipage}[b]{\linewidth}\raggedright
Score-5 mark
\end{minipage} \\
\midrule\noalign{}
\endhead
\bottomrule\noalign{}
\endlastfoot
\textbf{S1. Instructions} & Project/team/org rules attached to a
workspace & Versioned, two-loop, Auto-gated promotion \\
\textbf{S2. Skills} & Procedural artifacts (code/text) loaded on demand
& Skill bank self-modifies on failure, unit-test gated, versioned,
promoted \\
\textbf{S3. Memory} & Cross-session experience (episodic / semantic /
procedural) & Provenance + versioning + correction pathways \\
\textbf{S4. Tools} & Vetted action surface + context schemas &
Domain-curated bundle is the product; eval-gated tool versions \\
\textbf{S5. Orchestration} & Multi-role graph + routing &
Population/policy updates from outer-loop signal \\
\textbf{S6. Governance} & Versioning, promotion, audit, rollback &
dev→staging→prod with eval thresholds + audit log + rollback \\
\end{longtable}
}

The 0--5 rubric is uniform across substrates: \textbf{0} ephemeral;
\textbf{1} persistent local; \textbf{2} persistent + tool attachment, no
feedback; \textbf{3} multi-scope or feedback but no automated promotion;
\textbf{4} automated feedback updates the scaffold (one loop closed);
\textbf{5} two-loop versioned promotion plus governance (both loops
closed). We adopt this six-coordinate decomposition because each
coordinate is independently editable in production systems; alternative
partitions (e.g., separating evals from governance) are possible, and
the descent argument of §3 is invariant to the partition so long as the
independent-editability property holds.

\subsection{5. Survey method}\label{survey-method}

\textbf{Scope and corpora.} The survey covers approximately 130 unique
LLM systems published or productionised between January 2023 and May
2026. The candidate pool was assembled by searching the public discourse
for systems described as ``self-improving,'' ``scaffolded,'' or
``agentic,'' then filtering by whether the system wraps a foundation
model with at least one persistent artifact updated from deployment
signal. Approximately 90 systems satisfy that \emph{patch-local
discriminator} (score ≥ 3 on at least one of the six substrates) and
form the \textbf{core corpus}. The remaining \textasciitilde38 systems
carry the name but fail the discriminator; we keep them as a
\textbf{contrast corpus} so the boundary is auditable rather than
rhetorical (Appendix D lists representative entries).

\textbf{Discriminator and composite criterion.} The score ≥ 3 cut admits
a system as patch-locally adaptive. A separate, stricter audit applied
to 22 candidate systems uses a \emph{composite two-loop} criterion: Loop
1 modifies a persisted artifact from session signal, \emph{and} Loop 2
has an explicit cross-context promotion mechanism with versioning or
lineage. Thirteen research systems pass the strict audit. Two industry
exemplars (Sierra Agent OS 2.0, Cognition Devin) are also evaluated
using vendor documentation as primary evidence; they are flagged in the
§8 matrix because the evidence tier is weaker. The strict count moves
under two natural relaxations of the criterion (counting versioning OR
lineage OR rollback as governance; admitting optimiser-state updates as
Loop-1 persistence), but the qualitative findings of §§6--9 survive
both. Appendix B gives the derivation.

\textbf{Evidence policy.} Each system carries one or more of: (T1)
peer-reviewed publication; (T2) arXiv preprint; (T3) production claim
with metrics; (T4) open-source artifact; (T5) industry-blog or
interview-level documentation. T1 and T4 carry the highest weight in
cluster analysis. T3 is admissible when the claim is concrete (specific
benchmark, named integration target). T5 is admissible only when
corroborated by a higher tier elsewhere. Tier flags propagate to cluster
claims through the per-row evidence column in Appendix A; production
claims are not pooled with peer-reviewed evidence when computing
cluster-level statistics.

\subsection{6. Substrate maturity is
uneven}\label{substrate-maturity-is-uneven}

Aggregating the core corpus by substrate gives the picture below. The
column \(n\) counts surveyed systems on that substrate, \emph{Mean} is
the average score, and \emph{Score-5 systems} are the few that close
both an automated inner loop and a governed two-loop promotion at the
substrate level.

{\def\LTcaptype{none} % do not increment counter
\begin{longtable}[]{@{}
  >{\raggedright\arraybackslash}p{(\linewidth - 8\tabcolsep) * \real{0.1400}}
  >{\raggedleft\arraybackslash}p{(\linewidth - 8\tabcolsep) * \real{0.0400}}
  >{\raggedleft\arraybackslash}p{(\linewidth - 8\tabcolsep) * \real{0.0500}}
  >{\raggedright\arraybackslash}p{(\linewidth - 8\tabcolsep) * \real{0.2200}}
  >{\raggedright\arraybackslash}p{(\linewidth - 8\tabcolsep) * \real{0.5500}}@{}}
\toprule\noalign{}
\begin{minipage}[b]{\linewidth}\raggedright
Substrate
\end{minipage} & \begin{minipage}[b]{\linewidth}\raggedleft
\(n\)
\end{minipage} & \begin{minipage}[b]{\linewidth}\raggedleft
Mean
\end{minipage} & \begin{minipage}[b]{\linewidth}\raggedright
Score-5 systems
\end{minipage} & \begin{minipage}[b]{\linewidth}\raggedright
Survey finding
\end{minipage} \\
\midrule\noalign{}
\endhead
\bottomrule\noalign{}
\endlastfoot
\textbf{S1. Instructions} & 12 & 2.83 & \emph{(none)} & Rules persist
across IDEs and agents (Claude Code's CLAUDE.md, Cursor Rules,
AGENTS.md, Windsurf, Continue.dev) but no surveyed instruction system
closes both an automated inner loop and a governed two-loop promotion.
Ceiling held at 4 by Claude Code and Replit \\
\textbf{S2. Skills} & 16 & 2.62 & SAGE, COSPLAY & Research frontier is
active (failure-triggered skill update); production governance is
weaker. Anthropic's skills repository is one-pool with PR review but no
automated eval gate \\
\textbf{S3. Memory} & 19 & 3.16 & Cuadros et al.~governed collaborative
memory & Cross-session memory is widespread; \emph{correctable,
versioned} memory is rare. Production memory systems (ChatGPT Memory,
Claude Memory, Mem0, Zep, MemGPT) treat memory as append-only \\
\textbf{S4. Tools} & 19 & 3.63 & MCP, Harvey, Hippocratic AI & Most
mature production substrate. Tool bundles are the commercial unit; MCP
is the cross-vendor standard with broad public-server ecosystem and
substantial enterprise adoption \\
\textbf{S5. Orchestration} & 19 & 2.68 & \emph{(none)} & Multi-agent
topologies exist (LangGraph, AutoGen, CrewAI) but topology
\emph{evolution} is rare; MAE comes closest with RL co-evolution of
Proposer/Solver/Judge population but lacks Loop-2 cross-context
promotion \\
\textbf{S6. Governance} & 21 & 3.05 & Braintrust, Vellum, LangSmith Hub,
AGENTS.md/AAIF & High production maturity for prompts-as-code: immutable
commits, eval-gated PR merge, sub-five-minute rollback, typically
without inner-loop adaptation \\
\textbf{INTEGRATED} & 34 & 3.74 & NanoResearch, AutoAgent, SkillRL & The
high-water mark when present, but the cluster is dominated by research
systems lacking enterprise governance \\
\end{longtable}
}

Two things stand out. First, the field has independently converged on
the six substrates: every mature system has a recognizable instance of
each, and the names differ across vendors but the role does not. Second,
no system matures all six at once. Research systems learn fast and lack
enterprise governance. Production systems govern well and learn slowly.
Open standards solve artifact portability but not adaptation.
Vertical-bundle vendors solve patch specificity but rarely expose
cross-tenant promotion. The maturity gradient is not a quirk of the
corpus; it tracks which constraints each kind of system was built to
satisfy.

Five absences fall out of the discriminator and recur as Paper-4 targets
in §11. Failure-triggered memory is rare in production (most deployed
memory systems record preferences and successful facts, not residuals).
Cross-customer skill promotion is absent (Anthropic's skills repository
is one-pool; AGENTS.md is cross-vendor but not cross-customer). Memory
versioning is mostly unsolved outside the Cuadros et al.~framework.
Scaffold-level curriculum (which sessions inform the gradient) is
unspecified everywhere; every session contributes equally. Held-out
evals at Loop-2 promotion are absent, which is the missing piece that
would catch a scaffold over-fitting one deployment before it propagates
to another.

\subsection{7. Four clusters}\label{four-clusters}

Four archetypal clusters recur across the core corpus. They differ in
where on the \emph{patchwork to descent} gradient each cluster sits and
in which loop is automated. This section cuts the corpus by architectural
archetype; §8 re-cuts the composite-two-loop subset by gate type.

{\def\LTcaptype{none} % do not increment counter
\begin{longtable}[]{@{}
  >{\raggedright\arraybackslash}p{(\linewidth - 6\tabcolsep) * \real{0.2500}}
  >{\raggedright\arraybackslash}p{(\linewidth - 6\tabcolsep) * \real{0.2500}}
  >{\raggedright\arraybackslash}p{(\linewidth - 6\tabcolsep) * \real{0.2500}}
  >{\raggedright\arraybackslash}p{(\linewidth - 6\tabcolsep) * \real{0.2500}}@{}}
\toprule\noalign{}
\begin{minipage}[b]{\linewidth}\raggedright
Cluster
\end{minipage} & \begin{minipage}[b]{\linewidth}\raggedright
Representative systems (year, PP score)
\end{minipage} & \begin{minipage}[b]{\linewidth}\raggedright
Loop 1 / Loop 2
\end{minipage} & \begin{minipage}[b]{\linewidth}\raggedright
What the cluster proves
\end{minipage} \\
\midrule\noalign{}
\endhead
\bottomrule\noalign{}
\endlastfoot
Research full-auto & NanoResearch (2026, 5), SkillOpt (2026, 5), SkillRL
(2026, 5), AlphaEvolve (2025, 4), DGM (2025, 4), ADAS (2024, 4), Voyager
(2023, 4) & Auto / Auto & Both loops can be fully automated under
algorithmic gates; benchmark gains compound; enterprise governance is
uniformly absent \\
Production hybrid & SkillForge (2026, 4), CASCADE (2025, 4.5), Sierra OS
2.0 (2025, 3)\footnote{Sierra Agent OS 2.0 also fits the \emph{open
  standards / vertical-bundle} discussion below because its commercial
  product wraps a vertical workflow around the standardised agent-OS
  surface; we place it primarily in \emph{Production hybrid} on the
  basis of its Loop-1 and Loop-2 gate types.}, Cognition Devin (2024, 3)
& Auto / Human & Algorithmic Loop-1 gate plus reviewer-judgment Loop-2
gate is the production-viable cell; SkillForge is the architectural
exemplar \\
Governance-first & Braintrust, Vellum, LangSmith Hub (all 2024, S6 = 5)
& None / Auto & Eval-gated promotion exists; candidate generation is
manual; no patch-local inner loop \\
Open standards & AGENTS.md (2025), MCP (2024), Claude Skills (2025),
Cursor Rules (2024) & Mixed / Mixed & Vendor-cross convergence on
artifact format and tool attachment; tens of thousands of AGENTS.md
repositories; broad MCP server ecosystem with substantial enterprise
adoption \\
\end{longtable}
}

\emph{Research full-auto} systems demonstrate that the architecture
works under fully algorithmic gates. NanoResearch co-evolves Skills,
Memory, and Policy with a semantic-merge orchestrator. SkillOpt
operationalises the deep-learning-optimizer analogy explicitly, with
held-out validation gating, rejected-edit buffers, cosine-decay textual
learning rate, and epoch-wise slow/meta consolidation that together
resemble standard SGD machinery applied to a markdown artifact (§3.5;
Appendix G). SkillRL uses RL reward distillation with a hard threshold
at task-category accuracy below 0.4. AlphaEvolve runs a Gemini-driven
ensemble with a cascade evaluator and a MAP-Elites archive, deployed at
Google to gain +0.7\% Borg compute worldwide, +23\% on the Gemini
training kernel, +32.5\% on FlashAttention, and the first general-field
improvement over Strassen's bound on 4×4 matrix multiplication. DGM
modifies its own source code with archive-based parent selection and
lifts SWE-bench from 20\% to 50\%. The cluster maximises velocity. Enterprise governance
is uniformly absent.

\emph{Production hybrid} systems pair an algorithmic Loop-1 gate with a
reviewer-judgment Loop-2 gate. SkillForge drives auto-generated skill
diffs through an LLM-judge with \textgreater90\% human agreement (Loop
1), then routes what the judge passes through human support engineers
(Loop 2), with VFS version commits providing lineage. CASCADE reports
the largest within-benchmark ablation gain in the corpus (+57.9 pp on
SciSkillBench), driven by memory consolidation gated by human-agent
collaboration; the headline gain is not comparable across benchmarks
because SciSkillBench was designed to reward the very mechanism CASCADE
adds. Sierra Agent OS 2.0 routes AI-driven Insights
through human-authored Expert Answers in GitHub-style Workspaces.
Cognition Devin ships dynamic in-session tool creation with a 67\%
PR-merge rate, where the merged PRs are the Loop-2 promotion. This is
the production-viable cell.

\emph{Governance-first} systems automate Loop 2 alone. Braintrust,
Vellum, LangSmith Hub, W\&B Weave, Agenta, and LangFuse all instantiate
the \emph{prompts-as-code} pattern: immutable commits, semantic-version
tag pointers, PR-blocked merges on eval regression, sub-five-minute
rollback. Loop 1 is absent because the candidate artifacts are
engineer-authored at their desks rather than triggered by session
feedback. We frame this as \emph{MLOps for scaffolds}, distinct from
MLOps for weights.

\emph{Open standards} systems converge across vendors on the artifact
format and the tool layer. AGENTS.md is now in over 60,000 GitHub
repositories with measured efficiency gains: \citet{lulla2026empirical}
report roughly -29\% wall-clock and -17\% tokens with no quality loss. MCP
has become the cross-vendor tool standard with thousands of public
servers and substantial enterprise adoption. Claude Agent Skills, Cursor
Rules, GitHub Copilot custom instructions, Windsurf rules, Continue.dev
rules, Zed AI rules, and Replit \texttt{replit.md} all converge on a
git-tracked markdown artifact attached to a workspace, with an MCP-style
tool layer for actions. Convergence across vendors is the strongest
evidence the architecture is real rather than any one vendor's local
optimum.

The four clusters do not partition the corpus (many systems sit between
two clusters), but they capture the architectural variation.
Vertical-bundle vendors (Harvey for legal, Hippocratic AI for
healthcare) sit between \emph{open standards} and \emph{production
hybrid}: the bundle is the commercial unit; the underlying model is
mostly the same one a competitor would use. One result from the
orchestration literature also constrains the picture.
\citet{tran2026singleagent} report that single-agent LLMs match or beat
multi-agent systems on multi-hop reasoning when thinking-token budgets
are matched. The implication is that observed maturity gaps on S5 may
reflect the rarity of governed topology \emph{evolution}, and the
conflation of agent-count with compute-budget, rather than a structural
deficit of multi-agent designs.

\subsection{8. The two-loop design
space}\label{the-two-loop-design-space}

The composite-two-loop systems sit in a 3×3 matrix indexed by gate type
on each loop.

{\def\LTcaptype{none} % do not increment counter
\begin{longtable}[]{@{}
  >{\raggedright\arraybackslash}p{(\linewidth - 6\tabcolsep) * \real{0.2500}}
  >{\raggedright\arraybackslash}p{(\linewidth - 6\tabcolsep) * \real{0.2500}}
  >{\raggedright\arraybackslash}p{(\linewidth - 6\tabcolsep) * \real{0.2500}}
  >{\raggedright\arraybackslash}p{(\linewidth - 6\tabcolsep) * \real{0.2500}}@{}}
\toprule\noalign{}
\begin{minipage}[b]{\linewidth}\raggedright
\end{minipage} & \begin{minipage}[b]{\linewidth}\raggedright
\textbf{Loop-2 Auto-gated}
\end{minipage} & \begin{minipage}[b]{\linewidth}\raggedright
\textbf{Loop-2 Human-gated}
\end{minipage} & \begin{minipage}[b]{\linewidth}\raggedright
\textbf{Loop-2 None}
\end{minipage} \\
\midrule\noalign{}
\endhead
\bottomrule\noalign{}
\endlastfoot
\textbf{Loop-1 Auto-gated} & NanoResearch, SkillOpt, SkillRL,
AlphaEvolve, ADAS, DGM, Voyager, SAGE, EvolveR, CoEvoSkills, AutoAgent,
AutoManual & SkillForge, CASCADE, Sierra OS 2.0†, Cognition Devin† &
MAE*, MetaGen*, MetaReflection*, CLIN* \\
\textbf{Loop-1 Human-gated} & \textbf{{[}empty by engineering logic{]}}
& \textbf{{[}empty in surveyed corpus{]}} & \emph{(none)} \\
\textbf{Loop-1 None} & ExpeL (weak), Braintrust*, LangSmith Hub*,
Vellum* & Anthropic skills repo*, DSPy*, TextGrad* & Self-Refine*,
ReAct*, Mem0* \\
\end{longtable}
}

\emph{Asterisks fail at least one composite-two-loop criterion; included
for contrast.} † Industry exemplar evaluated using vendor documentation
as primary evidence.

Three cells are populated. \textbf{{[}Auto × Auto{]}} holds eleven
research systems with machine-evaluated decision criteria on both loops.
\textbf{{[}Auto × Human{]}} is the production-viable cell: Loop 1 fires
at machine speed under an algorithmic criterion, Loop 2 requires human
approval before propagation. \textbf{{[}None × Auto{]}} holds the
governance-first prompts-as-code systems with eval-gated CI but no
failure-triggered candidate generation. Three cells are empty for
distinct structural reasons. \textbf{{[}Human × Auto{]}} is empty by
engineering logic: if humans author candidates manually, automated
promotion offers little leverage. \textbf{{[}Human × Human{]}} is empty
in the surveyed corpus; a candidate occupant would be a solo or
small-team workflow where the marginal cost of a bad auto-merge exceeds
the marginal cost of human latency, and no public example was found in
the May 2026 window. \textbf{{[}Human × None{]}} is structurally
improbable: a system requiring human authorship on Loop 1 without any
Loop 2 collapses to engineer-edits-a-config, and is not patch-local in
the sense of §3.

Two cross-cutting findings sharpen the picture: noise reduction via
batching is largely absent (every auto-gated system fires at \(n = 1\)
per generation step, with SkillOpt a partial exception), and multi-layer
hierarchies are under-explored (only AlphaEvolve shows explicit depth).
The full cell-by-cell discussion, these two findings in detail, the
closest-approaches enumeration that supports §9, and the DGM dual-mode
footnote are in Appendix F.

\subsection{9. The missing architecture: governed cross-tenant scaffold
optimization}\label{the-missing-architecture-governed-cross-tenant-scaffold-optimization}

Four properties define the architecture that no surveyed system fully
instantiates:

\begin{itemize}
\tightlist
\item
  \begin{enumerate}
  \def\labelenumi{(\alph{enumi})}
  \tightlist
  \item
    \textbf{Auto-gated Loop 1.} The commit criterion is an algorithmic
    threshold.
  \end{enumerate}
\item
  \begin{enumerate}
  \def\labelenumi{(\alph{enumi})}
  \setcounter{enumi}{1}
  \tightlist
  \item
    \textbf{Multi-tenant cross-org promotion.} Loop 2 routes artifacts
    between organisations.
  \end{enumerate}
\item
  \begin{enumerate}
  \def\labelenumi{(\alph{enumi})}
  \setcounter{enumi}{2}
  \tightlist
  \item
    \textbf{Versioned lineage with RBAC.} Promoted artifacts carry
    semantic versions, audit trails, and access-control policies.
  \end{enumerate}
\item
  \begin{enumerate}
  \def\labelenumi{(\alph{enumi})}
  \setcounter{enumi}{3}
  \tightlist
  \item
    \textbf{Rollback.} A bad promotion can be reverted without redeploy.
  \end{enumerate}
\end{itemize}

The closest approaches partition the requirements. AlphaEvolve has (a)
and partial (b) within a single organisation. SkillForge has (a) and (b)
within one enterprise with partial (c) via VFS commits. SkillOpt has (a)
fully and partial (c) via the exported \texttt{best\_skill.md} artifact
and the protected slow-update field, and reports positive cross-model /
cross-harness / cross-benchmark transfer as empirical evidence-radius
measurement; it does not implement (b) cross-organisational promotion or
(d) explicit rollback. CORAL \citep{qu2026coral} has partial (a) via
heartbeat-gated promotion and asynchronous multi-agent Loop 2 within one
system; no (b) cross-organisational scope, no formal (c) RBAC, no (d)
rollback. PACE \citep{ling2026pace} has (a) at the small-model
production timescale and a fast/slow two-timescale Loop-1 + Loop-2 that
consolidates across one deployment's own episodes; no cross-deployment
(b), no formal (c) or (d). Meta-Harness \citep{lee2026metaharness} has
(a) over the harness substrate (S5) with a filesystem-of-candidates
credit-assignment record; no Loop-2 cross-organisational promotion, no
RBAC or rollback. Sierra Agent OS 2.0 has (b) and partial (a). The
governance-first cluster has (c) and (d) explicitly but no (a) or (b) in
the strict sense. The convergent evidence sharpens rather than reduces
the §9 gap: every closest approach satisfies a non-trivial subset of (a)
and (c), but no member yet provides (b) or (d) under any governance
regime. The full enumeration is in Appendix F, and Appendix G locates
each convergent system.

In ML terms, what is missing is a privacy-preserving cross-tenant
aggregator with an automated eval gate at the scaffold layer. Such an
architecture would carry five parts: per-tenant artifact pools, each
tagged with where its artifacts came from; a way to aggregate across
tenants that the operator can tune for privacy; an automated eval gate
that blocks a promotion whenever it does not lower patch loss; versioned
lineage with RBAC; and a rollback envelope on every promoted artifact.
The DP-FedAvg framing from the federated-learning literature is a useful
analogy for this target, but only if artifacts are first embedded as
mergeable update representations. In surveyed practice, Loop 2 is
selection and review rather than parameter averaging.

Naming the missing architecture converts a vague gap into a
four-property audit criterion that any future system can be measured
against. Whether it \emph{should} be filled is deployment-dependent. In
safety-critical domains, reviewer-judgment gates may be a design feature
rather than a defect. But no public system currently fills it, and that
is what matters for the next two years of architecture work. Whether the
optimizer-equipped composite reaches deployment-level reliability is a
deployment-by-deployment empirical question; the paper's contribution is
to name what \emph{would have to be true} of the optimizer for that
question to even be testable.

\subsection{10. A new failure surface}\label{a-new-failure-surface}

Patchwork scaffold maintenance introduces failure modes weight-only
systems do not have. Each one is also a missing piece of the optimizer.

\emph{Scaffold bloat} is the dominant near-term risk. IFScale
\citep{jaroslawicz2025ifscale} reports latency growing roughly 35× as
instruction count scales from 10 to 250, and the best frontier model in
that benchmark drops to 68\% accuracy at 500 instructions, with weaker
models degrading further. The underlying mechanism {[}lost-in-the-middle,
instruction interference; \citet{liu2023lostmiddle}{]} persists across
model generations even when individual frontier-model scaling improves.
Ratchet \citep{zhang2026ratchet} supplies direct empirical evidence that
bloat is also a \emph{signal-quality} problem and not only a latency one,
and that lifecycle hygiene is a precondition for residual capture and
mode discovery rather than an alternative to descent (Appendix G.7). The
architectural mitigation is a complexity penalty \(\Omega(z)\) in the
gate, explicit pruning in the loop, and hygiene operations treated as
first-class.

\emph{Poisoned memory} is the dominant security risk. Because memory is
persistent and is read back into later prompts, an attacker who gets one
malicious entry written can have it retrieved and acted on long after
the attack, in sessions that never saw the original input. PoisonedRAG
reports 90\%+ attack success rate against standard RAG; Memory Control
Flow Attacks report \textgreater90\% vulnerability across major LLM
agent frameworks. The mitigation is architectural too: record where each
memory entry came from (provenance), gate what is allowed to be written,
and keep a rollback envelope on accepted entries so a poisoned one can
be removed.

\emph{Prompt injection via tools and MCP} is the dominant cross-vendor
attack surface. The mechanism is that a tool's output, or a connected
MCP server, can smuggle instructions into the model's context that the
model then follows as if the user had issued them. CVE-2025-54136
(``MCPoison'') demonstrated a public MCP-server compromise pathway. The
architectural mitigation is to scope each tool's permissions narrowly
and to put tool bindings through the same promotion gate as any other
scaffold edit, treating them as first-class scaffold coordinates rather
than ambient configuration that nobody reviews.

\emph{Eval fragility} is the dominant ecosystem risk. The Leaderboard
Illusion documents up to 112\% relative performance gains when some
competitors get more access to the benchmark's own data than others ---
a score that reflects exposure to the test rather than true capability.
Eval suites are themselves patch-local artifacts and must survive
Goodhart's Law: once a metric becomes the target, optimisers learn to
game it. \citet{ong2026harnesseval} extend the diagnosis to scaffold
optimisers themselves: evaluating harness optimisers solely by
target-agent performance gains cannot distinguish \emph{informed} Loop-1
updates from Stage-2 trial-and-error, leaving open whether reported
gains are descent (§2 Stage 3) or random search (§2 Stage 2). The
architectural mitigation is held-out evals scoped within evidence radius
rather than promoted as universal benchmarks, plus direct evaluation of
optimiser updates against their attribution claims rather than only
their downstream scores.

\emph{Coordination failures} are documented in MAST: 14 failure modes,
dominantly role-confusion and context-loss.
The architectural mitigation is treating orchestration topology as a
scaffold coordinate subject to the same gate and rollback discipline as
instructions and memories.

\emph{Patch overfitting} is structural, not incidental: a scaffold tuned
to a deployment over-fits it, every time. The architectural mitigation
is \emph{evidence radius} discipline at promotion --- do not propagate a
fix beyond the scope that validated it.

\emph{Plasticity loss at the scaffold level} is the slow-moving risk.
Context rot, instruction conflict, memory pollution, and tool-version
drift produce a scaffold-level analogue of weight-level plasticity loss.
No surveyed system treats pruning or expiration as a first-class
operation. The architectural mitigation is to make expiry and rollback
explicit operations in the loop rather than retrospective hygiene.

The failure surface is real, quantified by 2024--2026 work, and largely
unmitigated in production. Every entry on it names a piece of the
optimizer that current systems do not yet have.

\subsection{11. Conclusion}\label{conclusion}

The field has built an external adaptation layer for LLMs, but not the
optimizer that should govern it. If reliability is increasingly repaired
through persistent scaffold artifacts, then production AI needs a theory
and engineering discipline for the discrete, gated, auditable learning
that happens there, not just for the differentiable learning inside the
model. This paper is one step toward that discipline: a survey of where
production sits on the \emph{patchwork to random-search to descent}
gradient of §2, a formalisation of the Stage-3 limit under stated
assumptions, and a named missing architecture where the optimizer most
plainly does not yet exist.

The field did not wait for frontier weights to become perfectly
reliable. It built a second learning system around them --- made of
markdown files, skills, memories, tools, eval suites, and rollback
buttons --- and what changed in 2025 and 2026 is that the combination
became dense enough across vendors to read as a single external
adaptation surface, most of which is still maintained as patchwork. Once
that surface is named, the missing pieces become concrete engineering
targets; the most consequential is governed cross-tenant scaffold
optimization, and that no public system yet fills it is the survey's
headline finding. The frontier model generalises. The local scaffold
specialises. Reliability comes from governing that specialisation:
turning today's patchwork into tomorrow's optimizer.

\begin{center}\rule{0.5\linewidth}{0.5pt}\end{center}

\subsection{Appendix A. Representative survey
rows}\label{appendix-a.-representative-survey-rows}

The 30 rows below are a representative spot-check of the corpus,
covering all six substrates, both research and production, and all four
clusters of §7. \emph{Loop 1} and \emph{Loop 2} columns apply to systems
satisfying the composite-two-loop criterion; \emph{n/a} means the loop
is absent. \emph{PP} is the integrated patch-local-adaptation score
(0--5); \emph{Evidence} abbreviates the strongest available source tier
(T1 peer-reviewed; T2 arXiv; T3 production claim; T4 open-source
artifact; T5 industry blog).

\begingroup\small

{\def\LTcaptype{none} % do not increment counter
\begin{longtable}[]{@{}
  >{\raggedright\arraybackslash}p{(\linewidth - 14\tabcolsep) * \real{0.1212}}
  >{\raggedleft\arraybackslash}p{(\linewidth - 14\tabcolsep) * \real{0.0404}}
  >{\raggedright\arraybackslash}p{(\linewidth - 14\tabcolsep) * \real{0.1111}}
  >{\raggedright\arraybackslash}p{(\linewidth - 14\tabcolsep) * \real{0.0707}}
  >{\raggedright\arraybackslash}p{(\linewidth - 14\tabcolsep) * \real{0.0707}}
  >{\raggedleft\arraybackslash}p{(\linewidth - 14\tabcolsep) * \real{0.0404}}
  >{\raggedright\arraybackslash}p{(\linewidth - 14\tabcolsep) * \real{0.1818}}
  >{\raggedright\arraybackslash}p{(\linewidth - 14\tabcolsep) * \real{0.3636}}@{}}
\toprule\noalign{}
\begin{minipage}[b]{\linewidth}\raggedright
System
\end{minipage} & \begin{minipage}[b]{\linewidth}\raggedleft
Year
\end{minipage} & \begin{minipage}[b]{\linewidth}\raggedright
Substrate
\end{minipage} & \begin{minipage}[b]{\linewidth}\raggedright
Loop 1
\end{minipage} & \begin{minipage}[b]{\linewidth}\raggedright
Loop 2
\end{minipage} & \begin{minipage}[b]{\linewidth}\raggedleft
PP
\end{minipage} & \begin{minipage}[b]{\linewidth}\raggedright
Evidence
\end{minipage} & \begin{minipage}[b]{\linewidth}\raggedright
Why included
\end{minipage} \\
\midrule\noalign{}
\endhead
\bottomrule\noalign{}
\endlastfoot
NanoResearch & 2026 & Integrated & Auto & Auto & 5 & T2
(arXiv:2605.10813) & Tri-level Skills/Memory/Policy co-evolution;
cleanest full-auto exemplar \\
AutoAgent & 2026 & Integrated & Auto & Auto & 5 & T2 (arXiv:2603.09716)
& Dual-cycle Execution+Evolution + elastic memory orchestration \\
SkillRL & 2026 & Integrated & Auto & Auto & 5 & T2 (arXiv:2602.08234) &
Recursive skill-augmented RL with \(\text{SR}<0.4\) threshold gate \\
SkillOpt & 2026 & Integrated (S2-led) & Auto & Auto & 5 & T2
(arXiv:2605.23904) & Explicit deep-learning-optimizer analogy; 52/52
cells best or tied; +23.5 / +24.8 / +19.1 GPT-5.5 avg over
direct/Codex/Claude Code; cross-model + cross-harness + cross-benchmark
transfer \\
AlphaEvolve & 2025 & Integrated & Auto & Auto & 4 & T2/T3
(arXiv:2506.13131) & Production proof-point: Borg +0.7\%, Gemini kernel
+23\%, FlashAttention +32.5\% \\
DGM & 2025 & Integrated & Auto & Auto & 4 & T2 (arXiv:2505.22954) &
Recursive self-modifying code; SWE-bench 20\%→50\% \\
ADAS & 2024 & Integrated & Auto & Auto & 4 & T2 (arXiv:2408.08435) &
Meta-agent search; Turing-complete agent representation \\
SkillForge & 2026 & Integrated & Auto & Human & 4 & T2
(arXiv:2604.08618) & Production hybrid exemplar; LLM-judge
\textgreater90\% + VFS versioning \\
CASCADE & 2025 & Integrated & Auto & Human & 4.5 & T2 (arXiv:2512.23880)
& Largest within-benchmark ablation gain: +57.9 pp on SciSkillBench \\
EvolveR & 2025 & Integrated & Auto & Auto & 3 & T2 (arXiv:2510.16079) &
Experience-driven lifecycle; accepted at ICML 2026 \\
PACE & 2026 & Integrated (small-model) & Auto & Auto & 4 & T2
(arXiv:2605.23019) & Two-timescale Loop-1 + Loop-2 self-evolution under
small-frozen-model production constraints \\
CORAL & 2026 & Integrated (S5+S3 async) & Auto & Auto & 4 & T2
(arXiv:2604.01658) & Only async multi-agent Loop-2 mechanism in the
corpus; shared persistent memory + heartbeat-gated promotion \\
MOSS & 2026 & Integrated (S2+S5 source) & Auto & n/a & 4 & T2
(arXiv:2605.22794) & Source-level rewriting; extends evolution to code
routing / hook ordering / dispatch unreachable from the text layer \\
RCL & 2026 & Integrated (formalisation) & Auto & n/a & 4 & T2
(arXiv:2604.03189) & Independent formalisation of artifact-layer
descent: credit assignment + \(\Omega(z)\) + overfitting + plasticity in
context space \\
GEPA & 2025 & (cross-substrate prompt opt) & Auto & n/a & 3 & T2
(arXiv:2507.19457) & Reflective prompt evolution; beats GRPO +6 avg /
+20 max with 35× fewer rollouts; SkillOpt's immediate predecessor \\
Voyager & 2023 & Integrated & Auto & Auto & 4 & T2 (arXiv:2305.16291) &
Foundational skill-library result; 3.3× items, 15.3× tech-tree
milestones \\
AGENTS.md & 2025 & S1 Instructions & n/a & n/a & 2 & T3/T4 (agents.md) &
Cross-vendor standard; 60k+ repos; Lulla 2026 measures −28.6\%
runtime \\
Claude Code CLAUDE.md & 2025 & S1 Instructions & n/a & n/a & 4 & T3/T4
(docs.anthropic.com) & Four scoping levels including org-IT policy +
auto-memory layer; ceiling for S1 \\
Cursor Rules & 2024 & S1 Instructions & n/a & n/a & 3 & T3/T4
(\url{docs.cursor.com/rules}) & \texttt{.cursor/rules/*.mdc} multi-scope
+ MCP attach; git-tracked \\
Anthropic Agent Skills & 2025 & S2 Skills & n/a & Human & 3 & T3
(\url{anthropic.com/engineering}) & Progressive disclosure spec;
LangChain replication 29\%→95\% pass-rate \\
SAGE & 2025 & S2 Skills & Auto & Auto & 5 & T2 (arXiv:2512.17102) &
+8.9\% SGC, 26\% fewer steps, 59\% fewer tokens on AppWorld \\
COSPLAY & 2026 & S2 Skills & Auto & Auto & 5 & T2 (arXiv:2604.20987) &
Co-evolving decision agent + learnable skill bank; long-horizon tasks \\
Ratchet & 2026 & S2 Skills (lifecycle) & n/a & n/a & 3 & T2
(arXiv:2605.22148) & Hygiene-as-precondition: \(+0.0\) pp without
hygiene vs.~\(+16.2\) pp with four hygiene mechanisms; Stage-1
prerequisite for Stage-3 descent \\
Reflexion & 2023 & S3 Memory & Auto & n/a & 4 & T1 (NeurIPS 2023;
arXiv:2303.11366) & Canonical failure-triggered episodic memory; +8\%
HotpotQA \\
Generative Agents & 2023 & S3 Memory & Auto & n/a & 4 & T1 (UIST 2023;
arXiv:2304.03442) & Reflection + importance memory primitive \\
Cuadros et al.~governed collaborative memory & 2026 & S3 Memory & Auto &
Auto & 5 & T2 (arXiv:2605.04264) & Only surveyed memory system with full
provenance + versioning + correction \\
MCP & 2024 & S4 Tools & n/a & n/a & 5 & T3 (modelcontextprotocol.io) &
Cross-vendor standard; thousands of public servers; substantial
enterprise adoption \\
Harvey & 2026 & S4 Tools & n/a & n/a & 5 & T3 (harvey.ai) & 400K
queries/day; 18,000+ workflows; 200+ legal data sources \\
Hippocratic AI & 2026 & S4 Tools & n/a & n/a & 5 & T3
(hippocraticai.com) & \$3.5B valuation; 30\% readmission reduction;
360\% care-capacity boost \\
Meta-Harness & 2026 & S5 Orchestration & Auto & n/a & 4 & T2
(arXiv:2603.28052) & End-to-end harness optimisation;
filesystem-of-candidates serves as accumulated history \(H_t\) for
credit assignment \\
Multi-Agent Evolve (MAE) & 2025 & S5 Orchestration & Auto & n/a & 4 & T2
(arXiv:2510.23595) & RL co-evolution of Proposer/Solver/Judge
population; one loop closed, no Loop-2 \\
Cemri et al.~(AG2) & 2025 & S5 Orchestration & n/a & n/a & 3 & T2
(arXiv:2503.13657) & Specialisation +4.5 pp on GSM-Plus (\(p = 0.03\));
MAST 14 failure modes \\
Braintrust & 2024 & S6 Governance & n/a & Auto & 5 & T3 (braintrust.dev)
& Prompts-as-code: immutable commits, eval-gated PR merge, sub-5 min
rollback \\
Vellum & 2024 & S6 Governance & n/a & Auto & 5 & T3 (vellum.ai) & Same
governance pattern; production deployment with eval thresholds \\
LangSmith Hub & 2024 & S6 Governance & n/a & Auto & 5 & T3
(\url{smith.langchain.com}) & Prompt repository with environment
promotion and eval blocking \\
Self-Refine* & 2023 & (contrast) & n/a & n/a & 1 & T2 (arXiv:2303.17651)
& Fails the patch-local discriminator: in-context iteration only \\
Mem0* & 2024 & (contrast) & n/a & n/a & 2 & T3/T4 (mem0.ai) & Fails Loop
2: per-user memory with no cross-context promotion path \\
Tran \& Kiela* & 2026 & (contrast) & n/a & n/a & 1 & T2
(arXiv:2604.02460) & Single-agent baselines match or beat multi-agent
under matched thinking-token budgets; orthogonal finding constraining S5
claims \\
\end{longtable}
}

\endgroup

\emph{Asterisks mark contrast systems included for definitional
clarity.}

\begin{center}\rule{0.5\linewidth}{0.5pt}\end{center}

\subsection{Appendix B. Data appendix}\label{appendix-b.-data-appendix}

\textbf{Corpus size.} The survey scores approximately 130 unique LLM
systems across the six substrates. Each system may be scored on multiple
substrates; the per-substrate row count therefore exceeds the
unique-system count.

\textbf{Core / contrast split.} Approximately 90 unique systems satisfy
the patch-local discriminator (score ≥ 3 on at least one substrate) and
form the core corpus. The remaining \textasciitilde38 systems form the
contrast corpus: commonly described as agentic or self-improving in
public discourse but failing at least one discriminator. Appendix D
enumerates representative contrast systems.

\textbf{Inclusion criteria.} Any system that (a) wraps a foundation
model with at least one persistent artifact updated from deployment
signal, or (b) is named in the public discourse as a ``self-improving,''
``scaffolded,'' or ``agentic'' system. The contrast set (Self-Refine,
ReAct, plain Mem0, DSPy/TextGrad) is included so the discriminator has
something to mark against.

\textbf{Exclusion criteria.} Systems whose only update mechanism is
weight-level fine-tuning or RLHF: these are model-side, not
scaffold-side, and belong to the self-evolution literature this paper
differentiates from. Systems with no public evidence of deployment
(white papers without code, demos, or production claim).

\textbf{Evidence tiers.} Each system carries one or more of: (T1)
peer-reviewed publication; (T2) arXiv preprint; (T3) production claim
with metrics; (T4) open-source artifact; (T5) industry-blog or
interview-level documentation. T1/T4 carry highest weight in cluster
analysis; T3 is admissible where the claim is concrete (specific
benchmark, named integration target); T5 is admissible only when
corroborated by T1--T4 elsewhere. Tier flags are propagated to cluster
claims through the per-row evidence column in Appendix A; T3 production
claims are not pooled with T1/T2 evidence when computing cluster-level
statistics.

\textbf{Evidence-tier distribution.} Of the core corpus, approximately
one quarter carry T1 evidence, roughly half carry T2 evidence, and the
remainder carry T3 or T4 evidence as the primary tier.

\textbf{Survey window and dating.} January 2023 -- May 2026.
Preprint-only systems carry T2 evidence and are flagged in Appendix A. A
fraction of T2 work is recent enough that ablations may yet be revised;
cluster-level claims are robust to dropping any single T2 system.

\textbf{Scoring discipline.} Each system received a
substrate-by-substrate 0--5 score against the rubric of §4 and an
integrated patch-local-adaptation score (PP) in \{0,1,2,3,3.5,4,4.5,5\}.
Half-scores were used sparingly when one criterion was clearly met and
another partially. Ambiguous cases were resolved conservatively
(downward). Scoring is single-rater; Appendix A and Appendix D together
expose representative per-system scores and exclusion reasons so readers
can audit ambiguous classifications independently.

\textbf{Composite two-loop criterion (sensitivity analysis).} The strict
count of thirteen research systems satisfying the composite-two-loop
criterion changes under two natural relaxations. Relaxation (a) counts
any of versioning, lineage, or rollback as Loop-2 governance; the count
rises to approximately seventeen by re-admitting MAE, MetaGen,
MetaReflection, and CLIN. Relaxation (b) admits session signal that
updates optimizer state without persisting an artifact distinct from
optimiser state; the count rises to approximately fifteen by
re-admitting DSPy and TextGrad. The qualitative findings (Auto×Auto
well-populated; {[}Human × Human{]} and {[}Human × Auto{]} empty in the
surveyed corpus; governance-first cluster in {[}None × Auto{]}; the
cross-tenant governed-optimisation cell empty) survive both relaxations.

\begin{center}\rule{0.5\linewidth}{0.5pt}\end{center}

\subsection{Appendix C. Gate type → estimator-family
mapping}\label{appendix-c.-gate-type-estimator-family-mapping}

The §3 formalism collapses production gates into a single per-edit gate
\(G_t\) that estimates the directional loss change
\(\widehat{\Delta}_{z_t} L_D\). In practice, surveyed gates instantiate
that abstraction in several distinct ways, and some are better read as
\emph{proposal-pruning} steps that precede the per-edit gate rather than
as direct \(\widehat{\Delta}\) estimators.

{\def\LTcaptype{none} % do not increment counter
\begin{longtable}[]{@{}
  >{\raggedright\arraybackslash}p{(\linewidth - 6\tabcolsep) * \real{0.2500}}
  >{\raggedright\arraybackslash}p{(\linewidth - 6\tabcolsep) * \real{0.2500}}
  >{\raggedright\arraybackslash}p{(\linewidth - 6\tabcolsep) * \real{0.2500}}
  >{\raggedright\arraybackslash}p{(\linewidth - 6\tabcolsep) * \real{0.2500}}@{}}
\toprule\noalign{}
\begin{minipage}[b]{\linewidth}\raggedright
Gate family
\end{minipage} & \begin{minipage}[b]{\linewidth}\raggedright
What it estimates
\end{minipage} & \begin{minipage}[b]{\linewidth}\raggedright
Slot in §3 formalism
\end{minipage} & \begin{minipage}[b]{\linewidth}\raggedright
Representative systems
\end{minipage} \\
\midrule\noalign{}
\endhead
\bottomrule\noalign{}
\endlastfoot
\textbf{Eval-cascade} & \(\widehat{\Delta} L_D\) on a held-out task
suite; cascaded through cheap-to-expensive evaluators & Per-edit gate
\(G_t\) & AlphaEvolve, Braintrust, Vellum, LangSmith Hub \\
\textbf{RL-threshold} & Reward improvement vs.~a deployment threshold &
Per-edit gate \(G_t\) & SkillRL, SAGE, EvolveR \\
\textbf{LLM-judge with calibration} & Pairwise- or rubric-scored
quality, calibrated against human agreement & Per-edit gate \(G_t\) &
SkillForge (Loop 1), CASCADE (consolidation gate) \\
\textbf{Surrogate verifier} & Symbolic, learned, or test-case verifier
producing pass/fail or graded score & Per-edit gate \(G_t\) & DGM
(sandboxed code), Voyager (in-game verifier) \\
\textbf{Reviewer-judgment} & Human credit assignment via PR review or
expert sign-off & Per-edit gate \(G_t\) (human-instantiated) &
SkillForge (Loop 2), Sierra Agent OS 2.0, Anthropic skills repo \\
\textbf{Semantic-merge} & Whether a candidate is semantically duplicated
by, or compatible with, existing scaffold artifacts &
\emph{Proposal-pruning} preceding \(G_t\) & NanoResearch (orchestrator
semantic-merge step) \\
\textbf{MAP-Elites / archive selection} & Whether a candidate populates
a previously unfilled cell in a behavior-diversity archive &
\emph{Proposal-pruning} preceding \(G_t\) & AlphaEvolve (archive), ADAS
(meta-agent search) \\
\textbf{Multi-sample agreement} & Whether \(k\) independent proposals or
evaluations converge on the same direction; false-accept rate decays as
\(p^k\) rather than variance shrinking as \(1/k\) & \emph{Aggregation
filter} (qualitatively distinct from mini-batch averaging) &
Under-explored in the surveyed corpus \\
\end{longtable}
}

The eval-cascade, RL-threshold, LLM-judge, surrogate-verifier, and
reviewer-judgment families directly populate the per-edit gate \(G_t\)
of §3 and inherit its descent guarantees (under the assumptions of
Appendix E). The semantic-merge and archive-selection families are not
estimators of \(\widehat{\Delta} L_D\). They constrain the proposal
distribution rather than judging individual candidate quality, and the
descent argument therefore applies to the composition (proposal-pruning
∘ per-edit gate). Multi-sample agreement filters are listed for
completeness because they are the correct ML analogue for ``wait for
\(k\) pieces of evidence before accepting an edit,'' which is \emph{not}
the same operation as mini-batch averaging and should not be conflated
with it.

\begin{center}\rule{0.5\linewidth}{0.5pt}\end{center}

\subsection{Appendix D. Contrast set}\label{appendix-d.-contrast-set}

The contrast corpus of approximately 38 systems is included to make the
patch-local discriminator audit-able. Representative entries:

{\def\LTcaptype{none} % do not increment counter
\begin{longtable}[]{@{}
  >{\raggedright\arraybackslash}p{(\linewidth - 6\tabcolsep) * \real{0.2500}}
  >{\raggedright\arraybackslash}p{(\linewidth - 6\tabcolsep) * \real{0.2500}}
  >{\raggedright\arraybackslash}p{(\linewidth - 6\tabcolsep) * \real{0.2500}}
  >{\raggedright\arraybackslash}p{(\linewidth - 6\tabcolsep) * \real{0.2500}}@{}}
\toprule\noalign{}
\begin{minipage}[b]{\linewidth}\raggedright
System
\end{minipage} & \begin{minipage}[b]{\linewidth}\raggedright
Year
\end{minipage} & \begin{minipage}[b]{\linewidth}\raggedright
Primary discriminator failure
\end{minipage} & \begin{minipage}[b]{\linewidth}\raggedright
Brief reason
\end{minipage} \\
\midrule\noalign{}
\endhead
\bottomrule\noalign{}
\endlastfoot
Self-Refine \citep{madaan2023selfrefine} & 2023 & No persistent artifact
& In-context iteration only \\
ReAct \citep{yao2023react} & 2023 & No persistent artifact &
Thought-act-observe loop is intra-session \\
Reflexion \citep{shinn2023reflexion} & 2023 & No Loop 2 & Per-agent
episodic memory without promotion \\
Generative Agents & 2023 & No Loop 2 & Per-agent reflection + importance
memory \\
ExpeL & 2023 & Loop boundary blurred & Experience bank built from
training tasks \\
Mem0 (vanilla) & 2024 & No Loop 2 & Per-user memory; no cross-context
promotion \\
Zep & 2024 & No Loop 2 & Per-agent/session knowledge graph \\
ChatGPT Memory & 2024 & No Loop 2 & Per-user persistence \\
Claude Memory & 2025 & No Loop 2 & Per-conversation-thread
persistence \\
MemGPT & 2023 & No Loop 2 & Virtual-context management per agent \\
DSPy & 2023 & No artifact distinct from optimiser state &
Prompt-template optimization; SkillOpt {[}\citet{yang2026skillopt};
§3.5; Appendix G{]} realises the artifact-bearing form. \\
TextGrad & 2024 & No artifact distinct from optimiser state & Same
boundary case as DSPy; SkillOpt {[}\citet{yang2026skillopt}; §3.5;
Appendix G{]} adds the persistent \texttt{best\_skill.md}, held-out
validation gate, rejected-edit memory, and textual learning-rate
schedule that TextGrad lacks. \\
MetaGen & 2026 & No Loop 2; no cross-session persistence & Role pool
dynamic within single inference call \\
MAE (Multi-Agent Evolve) & 2025 & No Loop 2 & Co-evolution within
Proposer/Solver/Judge triplet \\
MetaReflection & 2024 & No Loop 2 & Per-agent semantic memory \\
CLIN & 2023 & No cross-instance promotion & Causal abstractions
per-agent-per-environment \\
LangGraph, AutoGen, CrewAI & 2023--2024 & Framework only & Orchestration
frameworks; not deployed systems \\
Tran \& Kiela & 2026 & Single-agent baseline & Constrains S5 maturity
claims (see Appendix F) \\
W\&B Weave, Agenta, LangFuse & 2024 & No Loop 1 & Observability /
eval-as-code \\
Mobile-Agent-E, Cradle, OS-Copilot & 2024--2025 & No Loop 2 &
Self-evolving but per-session/per-deployment \\
Agent S, Agent S2 & 2024--2025 & No Loop 2 & Computer-use agentic
framework \\
ACE \citep{ace2025agentic} & 2025 & No Loop 2 in surveyed form &
Single-context evolving playbook \\
Cursor Rules (vanilla) & 2024 & No Loop 2 & Per-project artifacts \\
GitHub Copilot custom instructions, Windsurf, Continue.dev, Zed AI rules
& 2024 & No Loop 1 & Engineer-authored at repo base branch \\
Replit \texttt{replit.md} & 2025 & Loop 2 absent & Self-writing scaffold
but per-workspace \\
Aider conventions, Lovable knowledge & 2024 & No Loop 1 &
Engineer-authored convention files \\
\end{longtable}
}

The representative entries above cover the main exclusion families
(in-context-only; no Loop 2; framework-only; engineer-authored without
failure trigger; optimiser-state-only); finer-grained per-system
breakdowns are deferred to follow-on practitioner work.

\begin{center}\rule{0.5\linewidth}{0.5pt}\end{center}

\subsection{Appendix E. Formalization of artifact-layer
descent}\label{appendix-e.-formalization-of-artifact-layer-descent}

This appendix supplies the verification machinery for §3. The body
states the descent picture as a mechanism; this appendix derives the
one-step descent inequality, states the convergence proposition with its
assumptions, gives the full optimisation-analogy correspondence, walks
four concrete cases, and clarifies what credit assignment without chain
rule means.

\subsubsection{E.1 The mechanism stack}\label{e.1-the-mechanism-stack}

The abstract update rule of §3 expands into six operations that every
patch-local system must perform, whether explicitly engineered or
accidentally emergent.

\begin{enumerate}
\def\labelenumi{\arabic{enumi}.}
\tightlist
\item
  \textbf{Residual capture.} Log failures, corrections, rejected
  outputs, user edits, eval regressions, tool-call failures. The Loop-1
  signal source. Without it, \(z_t\) has no provenance.
\item
  \textbf{Mode discovery.} Cluster repeated failures into patch-specific
  error modes: the \emph{Architecture of Errors} catalogue, instantiated
  locally per patch.
\item
  \textbf{Delta synthesis.} Convert a recurring mode into a candidate
  scaffold edit: an instruction rewrite, a new skill, a memory
  correction, a tool binding, an orchestration change, an eval case.
  This produces \(z_t\).
\item
  \textbf{Delta validation.} Estimate whether the edit reduces \(L_D\)
  without collateral damage on adjacent inputs. The gate \(G_t\) is
  Auto-gated when the check is an algorithmic criterion, Human-gated
  when it is reviewer judgment.
\item
  \textbf{Promotion control.} Decide whether the accepted edit stays in
  the local fork or climbs the scope hierarchy: project, stack,
  organisation, cross-tenant. Loop 2.
\item
  \textbf{Drift and bloat control.} Expire stale rules, prune
  conflicting memories, roll back regressed skills, detect over-fitting.
  Without this, accumulated edits produce the scaffold-level analogue of
  plasticity loss.
\end{enumerate}

Each operation is the locus of a separate engineering discipline.
Surveyed systems implement them partially and unevenly: research
full-auto systems excel at (3) and (4); production governance leaders at
(5); industry hybrids at most of (1)--(5) but rarely (6). Implementation
patterns are deferred to follow-on practitioner work.

\subsubsection{E.2 Setup and notation}\label{e.2-setup-and-notation}

Let the frontier model have fixed parameters \(\theta_M\). Let the
scaffold be a structured parameter object
\[\theta_s \in \Theta_s = \Theta_1 \times \Theta_2 \times \Theta_3 \times \Theta_4 \times \Theta_5 \times \Theta_6,\]
where the six coordinates correspond to the six substrates of §4. The
space \(\Theta_s\) is discrete, symbolic, and versioned. For a
deployment patch \(D\), define patch loss as
\[L_D(\theta_s) = \mathbb{E}_{x \sim D}\!\left[\ell(\pi(\theta_M, \theta_s, x))\right].\]

A candidate edit \(z_t \in \mathcal{A}(\theta_s^t)\) is drawn from a
proposal process \(z_t \sim Q_t(z \mid H_t, \theta_s^t)\), where \(H_t\)
is the accumulated history of failures, corrections, rejected outputs,
eval regressions, reviewer comments, tool-call traces, and prior
accepted deltas, and conditioning on \(\theta_s^t\) ensures
admissibility at the current scaffold. \(Q_t\) is not uniform over
admissible edits: it concentrates on the failure modes with the largest
expected loss contribution. For a mode \(m\) that recurs with frequency
\(f_m\) and carries per-occurrence loss \(s_m\), that contribution is
\(f_m s_m\) --- repeatability times severity --- so high-frequency or
high-severity modes are proposed first. This is the proposal-side
counterpart to the gate's acceptance test, and what makes A3 (proposal
coverage), stated below, plausible rather than merely assumed. The
finite directional difference induced by the edit is
\[\Delta_z L_D(\theta_s) = L_D(\theta_s \oplus z) - L_D(\theta_s).\] A
useful edit is one with \(\Delta_z L_D < 0\). The gate's noisy estimate
of this quantity is \(\widehat{\Delta}_z L_D\).

\subsubsection{E.3 Descent inequality}\label{e.3-descent-inequality}

The gate operates on the \emph{estimator}, not the truth. The descent
argument distinguishes them.

\textbf{Assumption A1 (positive descent bias on accepted edits).} There
exists \(\gamma_t > 0\) such that
\[\mathbb{E}\!\left[\widehat{\Delta}_{z_t} L_D \mid G_t = 1, \theta_s^t\right] \leq -\gamma_t - \tau_t.\]
The gate's accept-threshold \(\tau_t\) ensures accepted edits clear the
margin under the estimator.

\textbf{Assumption A2 (bounded estimator bias on accepted edits).} There
exists \(\varepsilon_t \geq 0\) such that
\[\big|\mathbb{E}\!\left[\Delta_{z_t} L_D - \widehat{\Delta}_{z_t} L_D \mid G_t = 1, \theta_s^t\right]\big| \leq \varepsilon_t.\]

Combining A1 and A2 with the gated update rule of §3 and writing
\(p_t = \Pr(G_t = 1 \mid \theta_s^t)\),
\[\mathbb{E}\!\left[L_D(\theta_s^{t+1}) \mid \theta_s^t\right] \leq L_D(\theta_s^t) - p_t(\gamma_t + \tau_t) + p_t \varepsilon_t.\]
This is descent under a noisy estimator: per-step expected improvement
is the gate's margin minus its estimator bias, weighted by acceptance
probability.

\subsubsection{E.4 Convergence
proposition}\label{e.4-convergence-proposition}

\textbf{Proposition (convergence of loss, informal).} Assume A1, A2,
bounded loss \(0 \leq L_D \leq B\), and additionally:

\begin{itemize}
\tightlist
\item
  \textbf{A3 (proposal coverage).} At every non-\(G\)-stable scaffold
  there exists \(\delta > 0\) such that \(Q_t\) places probability at
  least \(q > 0\) on an admissible edit with
  \(\mathbb{E}[\Delta_z L_D] \leq -\delta\).
\item
  \textbf{A4 (effective compactness).} The reachable scaffold orbit
  takes values in a finite or precompact subset of \(\Theta_s\).
\item
  \textbf{A5 (persistent acceptance with summable bias).}
  \(\sum_t p_t (\gamma_t + \tau_t) = \infty\) and
  \(\sum_t p_t \varepsilon_t < \infty\).
\end{itemize}

Then \(L_D(\theta_s^t)\) is a non-negative supermartingale up to
summable noise; \(L_D(\theta_s^t) \to L_\infty\) almost surely for some
random \(L_\infty \geq 0\); and the limit set of \(\{\theta_s^t\}\) is
\textbf{\(G\)-stable}: no admissible edit clears the gate with negative
expected directional loss in the limit.

The proposition claims convergence of the \emph{loss} and stability of
the \emph{gate decision}, not convergence to a global or local optimum
of \(L_D\). A globally suboptimal \(G\)-stable scaffold is consistent
with the result: \(Q_t\) may never propose the improving edit, or
\(\widehat{\Delta}\) may be biased to reject an improving direction.
This is the ordinary target of local stochastic optimization under noisy
estimators rather than a guarantee of optimality.

\subsubsection{E.5 Optimisation-analogy
correspondence}\label{e.5-optimisation-analogy-correspondence}

{\def\LTcaptype{none} % do not increment counter
\begin{longtable}[]{@{}
  >{\raggedright\arraybackslash}p{(\linewidth - 2\tabcolsep) * \real{0.5000}}
  >{\raggedright\arraybackslash}p{(\linewidth - 2\tabcolsep) * \real{0.5000}}@{}}
\toprule\noalign{}
\begin{minipage}[b]{\linewidth}\raggedright
Optimization concept
\end{minipage} & \begin{minipage}[b]{\linewidth}\raggedright
Scaffold analogue
\end{minipage} \\
\midrule\noalign{}
\endhead
\bottomrule\noalign{}
\endlastfoot
Parameter vector & Instructions, skills, memories, tools, routing
graphs, eval gates, governance rules \\
Forward pass & Model--scaffold composite executing on patch input \\
Loss & Failure, user correction, eval regression, retrieval miss, tool
error, schema violation \\
Finite directional difference &
\(\Delta_z L_D(\theta_s) = L_D(\theta_s \oplus z) - L_D(\theta_s)\) \\
Gradient estimate & Credit assignment identifying the scaffold
coordinate responsible for the failure \\
SGD step & Accepted scaffold edit \\
Aggregated update & Multiple recurring failures aggregated before update
(either by sampling, \emph{mini-batch averaging}, or by
sign-of-direction agreement across \(k\) proposals, \emph{multi-sample
voting}) \\
Learning rate & Edit magnitude under \(\oplus\) \\
Update radius & Promotion scope: which loss is being descended on (USER,
PROJECT, STACK, TENANT, CORE) \\
Validation loss & Held-out patch eval before merge \\
Regularization & Complexity penalty \(\Omega(z)\), bloat control,
cross-patch regression checks \\
Overfitting & Scaffold improves on this patch and degrades elsewhere \\
Rollback & Reverting a harmful accepted update \\
\end{longtable}
}

\subsubsection{E.6 Credit assignment without chain
rule}\label{e.6-credit-assignment-without-chain-rule}

Scaffold systems do not compute chain-rule derivatives through markdown
files or tool registries, and the correspondence above deliberately
omits a ``backpropagation'' row. They do, however, perform \emph{credit
assignment} through a pipeline. If the final answer fails, the
responsible coordinate may be a retrieval configuration three layers
upstream, a missing tool schema constraint, an underspecified
instruction, a stale memory, or a bad orchestration edge. Engineers do
this manually; LLM-judges and surrogate verifiers do it
semi-automatically. The operation is \emph{attribution} (output error is
assigned to an upstream coordinate) rather than backpropagation. The
descent argument of §3 needs only that this attribution succeeds with
non-zero probability per failure (assumption (iv) of the scope box), not
that it is performed by chain rule.

\subsubsection{E.7 Worked cases}\label{e.7-worked-cases}

\textbf{RAG.} Suppose failures recur because retrieved context omits the
decisive clause. The observed loss is not ``the model hallucinated.''
Credit assignment traces the error to the retrieval coordinate:
chunking, query rewrite, metadata filter, embedding model, reranker,
citation policy, or context-packing rule. The candidate edit might
tighten chunk boundaries, add a metadata constraint, change the
reranker, or insert a regression eval containing the missed clause. The
gate tests whether retrieval recall and answer accuracy improve on the
patch eval without degrading adjacent cases. If accepted, the scaffold
has moved downhill on the local loss surface.

\textbf{Pitch review.} Suppose the system repeatedly overrates decks
claiming a ``huge market'' without numeric evidence. The failure cluster
identifies an under-specified rubric coordinate. The candidate edit
changes ``evaluate TAM size'' into ``require a numeric market estimate
and source; if absent, cap TAM score at 4/10.'' The gate runs the
revised rubric against held-out decks with human scores. If agreement
improves without damaging unrelated criteria, the edit is committed.

\textbf{Tool-using agent.} Suppose calls repeatedly fail because
optional parameters are underspecified. The responsible coordinate may
be the tool schema, the call policy, or the clarification rule. The
candidate edit adds required fields, constrained decoding, or a pre-call
uncertainty check. The gate estimates whether tool-call failure
decreases without increasing refusal, latency, or unnecessary
clarification.

\textbf{Memory.} Suppose a stale customer preference is repeatedly
retrieved and causes bad decisions. The failure is not a weight error.
The responsible coordinate is a memory entry plus its provenance,
priority, and expiry policy. The candidate edit corrects, expires, or
version-tags the memory. The gate checks whether the correction improves
future behavior without erasing useful history.

\subsubsection{E.8 Loop 1 and Loop 2 as descent steps of different
radius}\label{e.8-loop-1-and-loop-2-as-descent-steps-of-different-radius}

Loop 1 proposes and validates local descent steps on a single \(L_D\).
Loop 2 expands the update radius. It selects, merges, and promotes
accepted local edits so that they apply to a larger averaging set of
patch losses. USER-level acceptance descends on a single \(L_D\).
PROJECT promotion descends on a project-averaged loss. STACK promotion
widens the average further. CORE or cross-tenant promotion descends on a
federated average. The two-loop design space is therefore not just a
taxonomy of engineering patterns: it is a taxonomy of \emph{which loss
each surveyed system is descending on}, gated by which estimator.

\subsubsection{E.9 Overfitting in scaffold
adaptation}\label{e.9-overfitting-in-scaffold-adaptation}

In weight training, overfitting is usually a defect. In scaffold
adaptation, local overfitting is partly the point. Frontier weights are
trained to preserve cross-patch generality; they must smooth over local
conventions and rare workflow-specific failures. A deployment patch
needs the opposite: selective specialization to local residuals. The
scaffold is therefore an intentionally overfittable layer, but one whose
overfitting is scoped, inspectable, versioned, gated, and reversible. A
scaffold that improves one repository, hospital workflow, or legal
diligence process may degrade performance elsewhere. That is not a
contradiction; it is patch adaptation. The mistake is promoting that
edit beyond its evidence radius.

In non-stationary patches (\(D = D_t\)), the same descent mechanism
tracks a moving optimum rather than converging once. This is why recency
weighting, replay, pruning, expiry, and rollback are not optional
hygiene; they are the scaffold-level analogues of continual-learning
machinery.

\begin{center}\rule{0.5\linewidth}{0.5pt}\end{center}

\subsection{Appendix F. Survey detail}\label{appendix-f.-survey-detail}

This appendix carries the cell-by-cell discussion of the two-loop matrix
from §8, the closest-approaches enumeration that supports §9, and the
deployment-dependence argument. It supplies the verification material
for §§8--9.

\subsubsection{F.1 Two-loop matrix: cell-by-cell
discussion}\label{f.1-two-loop-matrix-cell-by-cell-discussion}

\textbf{{[}Auto × Auto{]}} is most populated (eleven systems). Every
member has a machine-evaluated decision criterion on both loops.
Velocity is the cluster's defining property; the lack of
cross-organisational governance is its uniform weakness.

\textbf{{[}Auto × Human{]}} is the production-viable cell: Loop 1 fires
at machine speed under an algorithmic criterion, Loop 2 requires human
approval before propagation. The pattern reconciles automation velocity
with reviewer accountability. SkillForge's LLM-judge + reviewer pipeline
is the architectural exemplar. CASCADE, Sierra, and Devin instantiate
variants; the industry exemplars (Sierra, Devin) are marked with †
because vendor documentation rather than peer-reviewed metrics is the
primary evidence.

\textbf{{[}None × Auto-gated{]}} holds the governance-first
prompts-as-code systems (Braintrust, LangSmith Hub, Vellum). They have a
real Loop 2 (eval-gated CI blocks promotion on regression), but no Loop
1 in the patch-local sense, because the candidate artifacts are authored
by engineers at their desks rather than triggered by session feedback.
The asterisks mark this: governance without inner-loop adaptation. The
cell is informative because it shows what survives when only the outer
loop is automated.

\textbf{{[}Human × Auto{]}} is empty by engineering logic: if humans
author candidates manually, automated promotion offers little leverage.
A speculative occupant would be ``expert curation + auto-distribution''
(expert-authored skills automatically gated through a held-out eval
before shipping); no published example exists.

\textbf{{[}Human × Human{]}} is empty in the surveyed corpus. A
candidate occupant would be a reviewer-authored-and-reviewer-promoted
artifact pipeline targeting solo and small-team workflows where
automation cost exceeds review latency. No public example was found
during the May 2026 survey window.

\textbf{{[}Human × None{]}} is structurally improbable: a system
requiring human authorship on Loop 1 but with no Loop 2 collapses to
``engineer-edits-a-config'' and is not patch-local in the sense of §3.

\textbf{DGM dual-mode note.} DGM operates in both {[}Auto × Auto{]}
(primary archive-based evolution) and {[}Auto × Human{]} (sandboxed
code-modification review) modes depending on deployment configuration;
we list it once in the dominant cell.

\subsubsection{F.2 Cross-cutting
findings}\label{f.2-cross-cutting-findings}

\textbf{Noise reduction via batching is largely absent in the surveyed
auto-gated systems.} Every auto-gated system fires at \(n = 1\) per
generation step, absorbing estimator noise statistically via the
per-edit gate rather than aggregating evidence across multiple
proposals. SkillOpt is a partial exception: it explicitly batches
evidence (rollout batch \(B\) and reflection minibatch \(B_m{=}8\)) and
decays edits via a learning-rate schedule, so the noisy-estimator
absorption is done by aggregation as well as by the per-edit gate; the
\(n{=}1\) characterisation still holds for the \emph{single accepted edit
per step}, but the proposal evidence is aggregated. Multi-sample
agreement filters (accept only when \(k\) independent proposals
converge) are a tractable design dimension that no surveyed system has
formalised. Their noise-reduction behaviour is qualitatively distinct
from mini-batch averaging (false-accept rate shrinks roughly as \(p^k\)
rather than variance shrinking as \(1/k\)), and the two should not be
conflated.

\textbf{Multi-layer hierarchies are under-explored.} Only AlphaEvolve
demonstrates explicit depth in the auto-gated cluster (via cascade
evaluator and MAP-Elites archive); the other twelve composite-two-loop
systems treat their artifact pool as flat.

\subsubsection{F.3 Closest approaches to governed cross-tenant scaffold
optimization}\label{f.3-closest-approaches-to-governed-cross-tenant-scaffold-optimization}

The four properties of governed cross-tenant scaffold optimization are
(a) auto-gated Loop 1, (b) multi-tenant cross-org promotion, (c)
versioned lineage with RBAC, (d) rollback. Closest approaches partition
the requirements:

\begin{itemize}
\tightlist
\item
  \textbf{AlphaEvolve} has (a) and partial (b) within a single
  organisation; no cross-org promotion mechanism.
\item
  \textbf{NanoResearch} has (a) at the research level; no enterprise
  governance layer described.
\item
  \textbf{SkillOpt} has (a) fully and partial (c) via the exported
  \texttt{best\_skill.md} artifact plus protected slow-update field;
  reports positive cross-model, cross-harness, and cross-benchmark
  transfer as empirical evidence-radius measurement; does not implement
  (b) cross-organisational promotion or (d) explicit rollback. The
  Appendix G walkthrough maps every SkillOpt design decision to its §3
  counterpart.
\item
  \textbf{CORAL} \citep{qu2026coral} has partial (a) via heartbeat-gated
  promotion and is the only system in the corpus that implements
  asynchronous multi-agent Loop 2 within one organisation through shared
  persistent memory; no (b) cross-organisational scope, no formal (c)
  RBAC, no (d) rollback.
\item
  \textbf{PACE} \citep{ling2026pace} has (a) at the small-model
  production timescale and a fast/slow two-timescale Loop 1 + Loop 2
  that consolidates across one deployment's own episodes rather than
  across deployments; no (b), no formal (c) or (d). PACE is the closest
  demonstration that the two-loop architecture is implementable under
  small-frozen-model constraints.
\item
  \textbf{Meta-Harness} \citep{lee2026metaharness} has (a) over the
  harness substrate (S5) with a filesystem-of-candidates
  credit-assignment record that functions as the accumulated history
  \(H_t\) of §3; no Loop-2 cross-organisational promotion, no RBAC, no
  rollback.
\item
  \textbf{SkillForge} has (a) and (b) within one enterprise, partial (c)
  via VFS commits; (d) is undocumented in the public report.
\item
  \textbf{Sierra Agent OS 2.0} has (b) and partial (a); the GitHub-style
  Workspaces provide some lineage but not RBAC-with-rollback in the
  strict sense.
\item
  \textbf{Vellum / Braintrust / LangSmith Hub} have (c) and (d)
  explicitly but no (a) or (b) in the strict sense; these are
  governance-first systems by construction.
\end{itemize}

No surveyed system combines (b), (c), and (d) under any gate type. The
cross-organisational governance triad is itself unfilled, independently
of whether Loop 1 is auto-gated.

\subsubsection{F.4 Deployment-dependence of the missing
architecture}\label{f.4-deployment-dependence-of-the-missing-architecture}

Whether the corner \emph{should} be filled is deployment-dependent. In
safety-critical domains (medical decision support, legal advice,
financial advice), reviewer-judgment gates may be a design feature
rather than a defect, and the human latency is the point. The narrower
survey claim is what matters: no public system currently combines
automated local scaffold evolution with governed cross-organisation
promotion, versioned lineage with RBAC, and rollback. Naming the corner
converts a vague gap into a specific four-property audit criterion that
any future system can be measured against. The shortest demonstrable
path appears to build on AlphaEvolve's cascade evaluator and MAP-Elites
archive, extending federation from within-DeepMind multi-target to
cross-customer multi-tenant. The architectural surface, not the
timeline, is what this paper claims.

\begin{center}\rule{0.5\linewidth}{0.5pt}\end{center}

\subsection{Appendix G. SkillOpt walkthrough and the convergent
cluster}\label{appendix-g.-skillopt-walkthrough-and-the-convergent-cluster}

This appendix expands §3.5 along two axes. The first six subsections
(G.1--G.3 on SkillOpt; G.4--G.8 on the convergent cluster) walk each
system mechanism-by-mechanism, stating its §3 mapping explicitly and
naming what it leaves open. The next subsection (G.9) synthesises four
claims supported by the cluster as a whole. The final two subsections
(G.10--G.11) carry the §9-property audit of SkillOpt and the published
cost-per-point figure.

\subsubsection{G.1 The convergent
landscape}\label{g.1-the-convergent-landscape}

SkillOpt \citep{yang2026skillopt} is the cleanest current instance of
artifact-layer descent in our survey, but it is one of at least fourteen
independent research threads in 2025--2026 that converge on the same
disciplined gradient-like update mechanism for the scaffold layer. The
cluster is striking because the convergence happens \emph{from different
starting points}: GEPA and Tiwari et al.~start from prompt-optimisation
and continual-learning respectively; RCL starts from cognitive-agent
learning; Meta-Harness, AHE, MOSS, CANTANTE, and Ong et al.~start from
harness engineering; MOCHA and Ratchet start from skill-deployment
constraints; ICT, Combee, CORAL, and PACE start from cross-task /
multi-agent / production-deployment angles. Each system instantiates a
different subset of §3's five conditions, and each exposes a different
gap.

The table below indexes the cluster by §3 mechanism and by what each
system leaves open. Subsequent subsections walk each cluster in detail.

{\def\LTcaptype{none} % do not increment counter
\begin{longtable}[]{@{}
  >{\raggedright\arraybackslash}p{(\linewidth - 8\tabcolsep) * \real{0.2000}}
  >{\raggedright\arraybackslash}p{(\linewidth - 8\tabcolsep) * \real{0.2000}}
  >{\raggedright\arraybackslash}p{(\linewidth - 8\tabcolsep) * \real{0.2000}}
  >{\raggedright\arraybackslash}p{(\linewidth - 8\tabcolsep) * \real{0.2000}}
  >{\raggedright\arraybackslash}p{(\linewidth - 8\tabcolsep) * \real{0.2000}}@{}}
\toprule\noalign{}
\begin{minipage}[b]{\linewidth}\raggedright
Cluster
\end{minipage} & \begin{minipage}[b]{\linewidth}\raggedright
Representative systems
\end{minipage} & \begin{minipage}[b]{\linewidth}\raggedright
What it instantiates in §3
\end{minipage} & \begin{minipage}[b]{\linewidth}\raggedright
Substrate (§4)
\end{minipage} & \begin{minipage}[b]{\linewidth}\raggedright
What it leaves open
\end{minipage} \\
\midrule\noalign{}
\endhead
\bottomrule\noalign{}
\endlastfoot
Reflective-evolution & GEPA \citep{agrawal2025gepa};
\citet{tiwari2026fastandslow} & \(Q_t\) over trajectory evidence with
Pareto-frontier proposal-pruning; two-timescale in-context / in-weights
motivation & (cross-substrate prompt opt) & persistent \(\theta_s\);
\(L_t\) budget; held-out \(G_t\); rejected-edit buffer \\
Reflective context & RCL \citep{vassilyev2026rcl} & \(\theta_s\) over
context space; generalisation penalty = \(\Omega(z)\); explicit naming
of credit-assignment / overfitting / plasticity-loss in context space &
S1+S2 (context) & cross-deployment Loop 2 \\
Harness optimisation & Meta-Harness \citep{lee2026metaharness}; AHE
\citep{lin2026ahe}; MOSS \citep{cai2026moss}; CANTANTE
\citep{zehle2026cantante}; \citet{ong2026harnesseval} & credit
assignment (§3 op 3); filesystem-of-candidates as \(H_t\); contrastive
multi-agent attribution; source-level extension beyond text-mutable & S5
(orchestration); S2+S5 (MOSS) & governed Loop 2 across orgs; RBAC;
rollback \\
Multi-objective + lifecycle & MOCHA \citep{tanjim2026mocha}; Ratchet
\citep{zhang2026ratchet} & Pareto frontier as the true shape of
\(\Omega(z)\); hygiene as precondition for residual capture and mode
discovery & S2 + deployment constraints & gradient-analogue closure;
gate strictness \\
Cross-task transfer & ICT \citep{shalev2026ict}; Combee
\citep{li2026combee}; CORAL \citep{qu2026coral}; PACE
\citep{ling2026pace} & empirical evidence-radius (Loop-2 scope)
measurement; async multi-agent Loop 2 within one org; two-timescale Loop
1 + Loop 2 under small-model constraints & Integrated; S3+S5 &
cross-organisational gate, lineage, rollback \\
Worked instance & SkillOpt \citep{yang2026skillopt} & all five §3
conditions: \(\theta_s\), \(L_D\), \(Q_t\)/\(z_t\), \(L_t\), \(G_t\) +
rejected-edit buffer + slow/meta update & S2-led Integrated & cross-org
Loop 2; RBAC; rollback \\
\end{longtable}
}

Read against §9, no member of the cluster combines (a) auto-gated Loop 1
with (b) cross-organisational Loop 2, (c) versioned lineage with RBAC,
and (d) rollback. The §3 picture is realised in fragments across the
cluster; the §9 gap is confirmed by every member.

\subsubsection{G.2 SkillOpt: full
correspondence}\label{g.2-skillopt-full-correspondence}

The table below extends the §3.5 compact correspondence with the §3 role
of each component and an algorithm-level pointer into SkillOpt's
published procedure (their Algorithm 1).

{\def\LTcaptype{none} % do not increment counter
\begin{longtable}[]{@{}
  >{\raggedright\arraybackslash}p{(\linewidth - 6\tabcolsep) * \real{0.2500}}
  >{\raggedright\arraybackslash}p{(\linewidth - 6\tabcolsep) * \real{0.2500}}
  >{\raggedright\arraybackslash}p{(\linewidth - 6\tabcolsep) * \real{0.2500}}
  >{\raggedright\arraybackslash}p{(\linewidth - 6\tabcolsep) * \real{0.2500}}@{}}
\toprule\noalign{}
\begin{minipage}[b]{\linewidth}\raggedright
§3 quantity
\end{minipage} & \begin{minipage}[b]{\linewidth}\raggedright
SkillOpt instantiation
\end{minipage} & \begin{minipage}[b]{\linewidth}\raggedright
§3 role explained
\end{minipage} & \begin{minipage}[b]{\linewidth}\raggedright
Algorithm pointer
\end{minipage} \\
\midrule\noalign{}
\endhead
\bottomrule\noalign{}
\endlastfoot
\(\theta_s\) (scaffold) & \texttt{best\_skill.md} --- a 379--1,995-token
markdown document & the trainable external state of the frozen target
model; persistent and exported as the deployment artifact &
\(s_{\text{best}}\) (Algorithm 1, lines 1, 22, 37) \\
\(L_D\) (patch loss) & benchmark-native hard score / exact-match on a
held-out test split & scalar quality signal that the optimiser is
implicitly minimising & \(\textsc{Evaluate}(M, h, s, D)\) (lines 2, 17,
37) \\
\(Q_t\) (proposal) & optimizer-model reflection over failure and success
minibatches; failures and successes analysed separately, then merged
with priority on failure fixes & proposal distribution conditioned on
accumulated rollout evidence and current scaffold & optimizer calls
\(O\) for failure analysis, success analysis, and merge (lines 9--11) \\
\(z_t\) (edit) & structured add/delete/replace patch on the skill (or
rewrite-from-suggestions in rewrite mode) & the candidate update applied
via \(\oplus\) to the current scaffold & rank/clip to \(L_t\) edits,
apply to obtain \(\tilde s\) (lines 12--13) \\
\(L_t\) (textual learning rate) & bounded per-step edit budget; constant
/ linear / cosine / autonomous schedules; default cosine with floor
\(L_t=2\) & edit-magnitude controller; the §3 \emph{learning rate} in
textual form & rank/clip step (line 12); schedule applied per epoch \\
\(G_t\) (gate) & held-out selection-split accept gate, strictly greater
than current selection score & the per-edit gate \(G_t\) --- accepts
only when \(\widehat{\Delta}_{z_t} L_D\) improves the selection score &
comparison at line 20; cache lookup at lines 14--18 \\
rejected-edit buffer & epoch-local store of failed proposals and the
score drop they caused & converts rejected steps into negative feedback
for later proposals within the same epoch; zero deployment cost &
\(\mathcal B\) updated at line 26; reset at line 5 \\
slow/meta update & epoch-end longitudinal guidance written to a
markup-fenced protected slow-update field, still passed through
\(D_{\text{sel}}\) & second-order consolidation across adjacent epochs;
the §3 ``step-level edits cannot overwrite the protected slow-update
field'' decomposition & lines 28--31; gate applied to injected
guidance \\
optimizer-side meta skill & \(m_{\text{meta}}\) --- a prompt-side
summary of which edits helped, which failed, and which failures
persisted & training-time guidance for \(Q_t\); never shipped with the
deployed artifact & lines 33--34 (training-side only) \\
\end{longtable}
}

The mapping is exact rather than analogical: every §3 quantity has a
named, code-level counterpart in SkillOpt's Algorithm 1, and every
SkillOpt control surface maps onto a §3 quantity. SkillOpt is therefore
the cleanest published instance we are aware of of the §3 picture as
\emph{mechanism} rather than \emph{metaphor} --- the five conditions of
§3 (measurable patch loss, residual capture, credit assignment,
candidate delta synthesis, directional-loss-change gate) are all
instantiated.

\subsubsection{G.3 What SkillOpt does that prior textual-optimisation
systems do
not}\label{g.3-what-skillopt-does-that-prior-textual-optimisation-systems-do-not}

SkillOpt's own experiment design compares it against six baselines under
an identical protocol (same target model, same held-out test split, same
scorer). The differentiation below frames each baseline in §3 vocabulary
and identifies the §3 mechanism the baseline lacks. The through-line is
the user-stated thesis: \emph{the scaffold/evolution layer needs a
disciplined gradient-like update mechanism}, and the live distinction
across baselines is which §3 conditions each one fails.

\begin{itemize}
\item
  \textbf{TextGrad} \citep{yuksekgonul2024textgrad}\textbf{.}
  Backpropagates textual feedback through \texttt{tg.Variable} objects
  in a compound system. In §3 terms: \(z_t\)-style edits exist, but
  \(\theta_s\) does not --- TextGrad has no artifact distinct from
  optimiser state. There is no \(L_t\) schedule, no held-out \(G_t\),
  and no rejected-edit buffer. SkillOpt's measured advantage over
  TextGrad on direct chat ranges from \(+5.9\) to \(+34.0\) points
  across benchmarks at GPT--5.5 scale and is uniformly positive at all
  seven target-model scales. This matches our Appendix D classification
  of TextGrad as failing the patch-local discriminator on persistence;
  SkillOpt closes precisely that gap.
\item
  \textbf{GEPA} \citep{agrawal2025gepa}\textbf{.} Reflective prompt
  evolution with Pareto-style multi-candidate selection. Closer than
  TextGrad in that it carries multiple candidates and a selection step.
  In §3 terms: \(Q_t\) is real, but the optimisation target is a prompt
  rather than a persistent skill artifact, \(G_t\) is not the
  strictly-greater held-out gate of §3, and there is no \(L_t\) bound
  --- accepted edits can wholesale rewrite. SkillOpt outperforms GEPA in
  the reported comparison; the per-cell deltas should be read directly
  from SkillOpt's results table rather than summarised here.
\item
  \textbf{Trace2Skill} \citep{ni2026trace2skill}\textbf{.} Distils
  trajectory-local lessons into reusable skill artifacts. In §3 terms:
  \(\theta_s\) exists, but there is no iterative \(G_t\) --- Trace2Skill
  mines skills offline rather than running a held-out validation loop
  that accepts or rejects each candidate. Our §3 mechanism stack
  therefore loses \emph{delta validation} (operation 4) and the
  \emph{rejected-edit buffer}. SkillOpt's measured gap over Trace2Skill
  is \(+5.0\) to \(+19.0\) points across benchmarks at GPT--5.5 direct
  chat.
\item
  \textbf{EvoSkill} \citep{alzubi2026evoskill}\textbf{.} Skill-folder
  evolution under failure analysis, deployed harness-side. The strongest
  baseline on the Codex and Claude Code harness rows. In §3 terms:
  \(\theta_s\) is a folder rather than a single artifact, and credit
  assignment via failure analysis is genuine, but there is no \(L_t\)
  bound on per-step folder edits and no rejected-edit buffer. SkillOpt
  beats EvoSkill by \(+14.0\) inside Codex and \(+3.2\) inside Claude
  Code on the GPT--5.5 five-benchmark average.
\item
  \textbf{Human skill} (expert-written, 145--516 tokens per benchmark).
  In §3 terms: \(\theta_s\) exists, but Loop 1 does not --- the artifact
  is static at deployment. Provides the no-Loop-1 ceiling.
\item
  \textbf{One-shot LLM skill} (generated zero-shot from a benchmark
  description by GPT--5.5, never updated). In §3 terms: \(\theta_s\)
  exists; Loop 1 does not; \(Q_t\) fires exactly once with no rollout
  evidence. The ``untrained scaffold'' baseline.
\end{itemize}

ABSTRAL \citep{song2026abstral} and EvoTest \citep{he2025evotest},
referenced in SkillOpt's related work but not measured in their table,
operate at the multi-agent design / test-time configuration layer rather
than the persistent skill layer; they are out of scope for the
artifact-layer-descent comparison.

Synthesising across all six measured baselines: the live difference is
whether the loop has (i) a persistent edit-target distinct from
optimiser state (\(\theta_s\)), (ii) a held-out validation gate on each
accept (\(G_t\) with \(\widehat\Delta\) estimator and accept threshold
\(\tau_t\)), (iii) negative-feedback retention (rejected-edit buffer
feeding \(Q_t\)), and (iv) a bounded learning rate with a schedule
(\(L_t\) + cosine/linear/constant/autonomous). SkillOpt is the first
published system in our survey that instantiates all four
simultaneously, and it does so as a single offline training loop whose
deployed artifact remains a 300--2,000-token markdown file requiring no
inference-time optimiser calls.

\subsubsection{G.4 GEPA, Tiwari, and the reflective-evolution
family}\label{g.4-gepa-tiwari-and-the-reflective-evolution-family}

The immediate predecessor to SkillOpt's optimiser framing is
\textbf{GEPA} \citep{agrawal2025gepa}. It shows that natural-language
reflection over execution trajectories can provide a richer adaptation
signal than sparse scalar rewards for prompt evolution in compound AI
systems. GEPA samples trajectories containing reasoning steps, tool
calls, and tool outputs; reflects on them to diagnose failures and
propose prompt updates; and maintains a Pareto frontier of candidate
prompts rather than greedily mutating only the current best. Across six
tasks, GEPA outperforms GRPO by 6\% on average and by up to 20\% while
using up to 35× fewer rollouts, and also outperforms MIPROv2 in the
reported aggregate comparison. In §3 vocabulary, GEPA strongly
instantiates the proposal side, \(Q_t\), through trajectory-conditioned
reflection, and uses Pareto-based candidate selection as a
proposal/pruning mechanism. What it does not instantiate is the full
SkillOpt-style artifact lifecycle: a bounded edit budget on a reusable
skill document, a strictly-greater held-out gate on each accepted skill
update, a rejected-edit buffer used as negative feedback, and an
exported deployment artifact with governance hooks.

\textbf{\citet{tiwari2026fastandslow}} supply the broader theoretical
framing from the optimisation side, arguing explicitly that restricting
learning to \emph{either} in-context (prompt optimisation) \emph{or}
in-weights (RL / fine-tuning) is unnecessary and costly. In-context
adaptation is cheap and rapid but capacity-bounded; in-weights
adaptation is expressive but induces catastrophic forgetting and loss of
plasticity \citep{dohare2024loss}. The scaffold layer of §2 is precisely
the substrate that absorbs what in-context adaptation can carry cheaply,
leaving the weights free to generalise. Tiwari et al.~supply the
two-timescale motivation from the optimisation side; §2 supplies the
matching motivation from the production-reliability side. The two
framings are independent and mutually reinforcing.

\subsubsection{G.5 Reflective Context Learning: a parallel
formalisation}\label{g.5-reflective-context-learning-a-parallel-formalisation}

The most direct independent convergence on the §3 formalisation comes
from \textbf{Reflective Context Learning} {[}RCL;
\citet{vassilyev2026rcl}{]}. Their opening observation is that ``the
fundamental problems of learning, including credit assignment,
overfitting, forgetting, local optima, and high-variance learning
signals, persist whether the learned object lies in parameter space or
context space,'' and that current methods are ``fragmented and ad hoc''
precisely because they address these challenges one at a time rather
than as a unified optimisation problem over the context. The parallel to
§3's argument --- that scaffold maintenance needs the discipline of
weight-space optimisation applied to an external artifact layer --- is
exact and arrived at independently.

RCL maps onto §3 as follows. The \emph{learned context} is \(\theta_s\).
The \emph{reflection signal} is a form of credit assignment identifying
which context coordinate caused a failure (§3 operation 3). The
\emph{generalisation penalty} that RCL introduces to prevent narrow
sample-specific rules is structurally the complexity penalty
\(\Omega(z)\) of the gate condition. Where RCL and §3 differ is in
Loop-2 scope: RCL is framed around single-agent cross-task
generalisation, while §3's Loop 2 is cross-deployment, cross-tenant, and
cross-organisational promotion. The convergence is nonetheless
significant: two research threads starting from different motivations
(production reliability in §3, cognitive agent learning in RCL) arrive
at the same formal observation that the classical learning problems do
not disappear when the learned object moves outside the weights.

\subsubsection{G.6 The harness-optimisation
cluster}\label{g.6-the-harness-optimisation-cluster}

A second convergent cluster focuses on the \emph{harness} --- the code
and configuration layer that wraps a frozen model --- rather than on
skill documents specifically. This cluster independently identifies
credit assignment as the central unsolved problem of scaffold evolution,
and its engineering solutions directly instantiate §3's operation 3
(mode discovery) and operation 4 (delta validation).

\textbf{Meta-Harness} \citep{lee2026metaharness} is the cleanest
representative. Its opening claim --- that ``the performance of large
language model systems depends not only on model weights, but also on
their harness: the code that determines what information to store,
retrieve, and present to the model'' --- is §1's scaffold-versus-weights
decomposition restated for the harness substrate specifically.
Meta-Harness then identifies that existing text optimisers ``compress
feedback too aggressively,'' losing the fine-grained signal needed for
credit assignment. Its solution is an outer-loop proposer with access to
source code, scores, and execution traces of all prior candidates
through a filesystem, so that the edit history itself becomes the
credit-assignment signal. In §3 terms, the filesystem of prior
candidates is a concrete form of the accumulated history \(H_t\) that
conditions the proposal distribution \(Q_t\), and the outer loop is a
Loop-1 closure over the harness substrate (§4, S5) rather than the skill
substrate (S2).

\textbf{Agentic Harness Engineering} {[}AHE; \citet{lin2026ahe}{]}
attacks the same problem from the observability side. The core diagnosis
is that harness engineering ``faces a heterogeneous action space across
editable components, voluminous trajectories that bury actionable
signal, and edits whose effect is hard to attribute.'' These three
failure modes map precisely onto §3's mechanism stack: a heterogeneous
action space is the coordinate-identification problem (op 2, mode
discovery); voluminous trajectories that bury signal is the
residual-capture problem (op 1); and effects that are hard to attribute
is operation 3, credit assignment. AHE's response --- three matched
observability pillars giving every editable harness component a
file-level representation --- is an engineering instantiation of the
coordinate decomposition that §3 assumes. The implicit claim is that
artifact-layer descent requires not just the optimisation loop but also
the instrumentation that makes scaffold coordinates individually
observable and attributable.

\textbf{MOSS} \citep{cai2026moss} extends the harness argument by
observing that all prior self-evolving agent systems ``confine evolution
to text-mutable artifacts --- skill files, prompt configurations, memory
schemas, workflow graphs --- and leave the agent harness untouched.''
Since routing, hook ordering, state invariants, and dispatch live in
code rather than in any text artifact, an entire class of structural
failure is, in MOSS's phrase, ``physically unreachable from the text
layer.'' MOSS therefore extends Loop-1 evolution to source-level code
rewriting. In §4 terms, MOSS evolves S2 (Skills) and S5 (Orchestration)
simultaneously, rather than treating the orchestration substrate as a
fixed wrapper around an evolvable skill layer. This matters for §9: a
governed cross-tenant scaffold optimiser that covers only text-mutable
artifacts leaves the code substrate ungoverned, and MOSS's diagnosis
suggests the gap is not merely inconvenient but structurally incomplete.

\textbf{CANTANTE} \citep{zehle2026cantante} attacks credit assignment
directly in the multi-agent setting. Its observation is that ``scores
are available only at the system level, whereas the parameters governing
agent behavior are local,'' making optimisation of multi-agent scaffolds
``fundamentally a credit-assignment problem.'' CANTANTE's solution is
contrastive: it decomposes system-level rewards into per-agent update
signals by contrasting rollouts of multiple jointly acting
configurations. In §3 vocabulary, CANTANTE provides a concrete mechanism
for operation 3 in the case where the scaffold is distributed across
agents rather than concentrated in a single artifact. The gap it closes
is the same one §3 identifies in Appendix E.6 --- that attribution
(output error assigned to an upstream scaffold coordinate) must succeed
with non-zero probability per failure --- but instantiated for the
multi-agent topology substrate (S5).

\textbf{Harness optimizer evaluation} \citep{ong2026harnesseval}
provides independent empirical evidence for a failure mode this paper
flags in §10: that evaluating scaffold optimisers ``solely by observing
target agents' performance gains'' misses intermediate erroneous edits
that hinder performance. Their finding --- that it is ``unclear whether
harness optimization is driven by optimizers' informed update actions or
simply trial-and-error'' --- is a direct empirical instantiation of the
distinction between Stage 2 (random artifact search) and Stage 3
(artifact-layer descent) from §2's table. The credit-assignment problem
is not merely theoretical: production harness optimisers that cannot
distinguish informed updates from lucky trial-and-error are operating in
Stage 2 even when they report benchmark improvements.

\subsubsection{G.7 Multi-objective and lifecycle
constraints}\label{g.7-multi-objective-and-lifecycle-constraints}

Two papers expose constraints on the §3 picture that single-loss descent
does not capture, both relevant to the §9 missing-architecture
discussion.

\textbf{MOCHA} \citep{tanjim2026mocha} observes that agent skills are
``multi-field artifacts subject to hard platform constraints:
description fields are truncated for routing, instruction bodies are
compacted via progressive disclosure, and co-resident skills compete for
limited context windows.'' Skill optimisation is therefore not a
single-loss descent problem but an inherently multi-objective one: a
skill must simultaneously maximise task performance and satisfy platform
limits, and optimising for task performance alone produces skills that
pass benchmarks but violate deployment constraints. MOCHA introduces
\emph{Chebyshev annealing} to navigate the constraint frontier. In §3
terms, the complexity penalty \(\Omega(z)\) in the gate condition is
precisely the mechanism MOCHA is operationalising, but §3's scalar
\(\Omega\) is a simplification: in production, the constraint surface is
a Pareto frontier, not a single penalty term. MOCHA's empirical
contribution is to show that ignoring this frontier --- treating skill
optimisation as unconstrained single-loss descent --- produces skills
that degrade on real deployment infrastructure even when benchmark
scores improve.

\textbf{Ratchet} \citep{zhang2026ratchet} makes an argument that is, in
some respects, orthogonal to the gradient-analogue framing and for that
reason worth engaging directly. Ratchet's opening finding --- that
LLM-authored skills deliver \(+0.0\) pp over no-skill baselines while
human-curated ones deliver \(+16.2\) pp --- suggests that the bottleneck
in self-evolving skill libraries is not the optimisation algorithm (the
mechanism for proposing and accepting edits) but \emph{lifecycle
management}: whether skills are written, retrieved, curated, and retired
correctly. Ratchet introduces four hygiene mechanisms --- outcome-driven
retirement, a bounded active-cap, meta-skill authoring, and pattern
canonicalisation --- and shows that these lift the \(+0.0\) pp floor
substantially before any gradient-analogue optimiser is applied.

This finding is not a refutation of the §3 picture but a clarification
of the precondition for it. Artifact-layer descent, as formalised in §3,
presupposes that the scaffold is maintained well enough for the gate
\(G_t\) to have a meaningful signal: if skill retrieval is so noisy that
the wrong skill is injected into every rollout, the proposal process
\(Q_t\) receives corrupted evidence and the descent guarantee degrades
to random search. Ratchet's hygiene mechanisms are, in §3 vocabulary,
prerequisites for operation 1 (residual capture) and operation 2 (mode
discovery) to function: outcome-driven retirement prevents stale
artifacts from polluting the residual signal; the bounded active-cap
prevents scaffold bloat from making retrieval unreliable (the §10
failure mode); pattern canonicalisation collapses duplicate and
near-duplicate skills so the proposal process conditions on a clean,
deduplicated artifact set rather than competing variants of the same
pattern. The practical implication is that the deployment path to
Stage-3 artifact-layer descent runs through Stage-1 hygiene as a
\emph{prerequisite}, not as an alternative. A system with strong hygiene
and no optimiser sits at Stage 1 with a clean scaffold; a system with a
strong optimiser and poor hygiene sits at Stage 2 with a noisy loss
signal. Stage 3 requires both.

\subsubsection{G.8 Cross-task transfer and evidence
radius}\label{g.8-cross-task-transfer-and-evidence-radius}

Several papers independently measure what §3 calls the \emph{evidence
radius} of a scaffold artifact --- how far an edit validated in one
context continues to improve performance when redeployed elsewhere ---
without using that vocabulary.

\textbf{Training Language Agents to Learn from Experience} {[}ICT;
\citet{shalev2026ict}{]} introduces the In-context Training framework
specifically to evaluate \emph{cross-task self-improvement}: whether
experience distilled from one set of tasks transfers to unseen tasks.
Their finding that ``current reflection-based methods can only
self-correct within a single task instance'' is a measurement of
evidence radius at its lower bound --- a radius of exactly one task ---
and their ICT framework is a systematic attempt to expand that radius
through supervised distillation of reflections into reusable lessons. In
§3 terms, ICT is an attempt to instrument Loop 2 empirically: to measure
whether a locally-validated edit (one that passed the gate on one task's
\(L_D\)) continues to improve performance when the promotion radius is
expanded to cover a distribution of unseen tasks.

\textbf{Combee} \citep{li2026combee} approaches the same measurement
from the scaling direction. ACE \citep{ace2025agentic} established that
context-as-living-artifact (a system prompt that accumulates, refines,
and organises strategies through agentic interaction) could improve over
static prompts on single-agent tasks. Combee extends this to
high-parallelism settings, arguing that ``existing methods primarily
focus on single-agent or low-parallelism settings,'' which
``fundamentally limits their ability to efficiently learn from a large
set of collected agentic traces.'' The parallel scaling Combee addresses
is a form of Loop-2 scope expansion: when many agents run
simultaneously, the trace pool available for the next proposal step is
drawn from a wider deployment distribution, and the promoted artifact is
being tested against a larger averaging set of patch losses. This is not
cross-organisational Loop-2 promotion, but it is a step toward it, and
Combee's empirical results measure how much of the evidence-radius
benefit can be captured by intra-organisation parallelism before the
cross-organisational barrier is reached.

\textbf{CORAL} \citep{qu2026coral} addresses the Loop-2 problem from the
multi-agent open-ended discovery angle. CORAL replaces rigid control
with long-running agents that explore, reflect, and collaborate through
shared persistent memory, asynchronous multi-agent execution, and
heartbeat-based interventions. The shared persistent memory is the
cross-agent promoted artifact: a finding validated in one agent's
exploration that is written to the shared store becomes available to all
agents, and the heartbeat mechanism provides a form of auto-gated Loop-2
promotion. CORAL does not implement the versioned lineage, RBAC, or
rollback that §9 requires, but it provides the only fully asynchronous
multi-agent Loop-2 promotion mechanism in the current survey corpus, and
its open-ended discovery framing makes the evidence-radius question
concrete.

\textbf{PACE} \citep{ling2026pace} closes this cluster with a
production-oriented two-timescale design for small language model
agents. The fast timescale adapts prompts, parsers, validators, and
control logic from failure signals within a deployment; the slow
timescale consolidates across episodes. The two-timescale structure maps
cleanly onto the Loop-1 / Loop-2 distinction of §3: Loop 1 fires at the
fast timescale, updating a single deployment's scaffold; Loop 2 fires at
the slow timescale, consolidating across the deployment's own episodes
rather than across deployments. PACE's contribution is to show that the
two-loop architecture is implementable under the resource constraints of
small, frozen production models --- not only under frontier-model
assumptions --- and that the slow/fast consolidation structure persists
even when the optimiser model and the target model are the same frozen
small model.

\subsubsection{G.9 Synthesis: what the convergence
establishes}\label{g.9-synthesis-what-the-convergence-establishes}

Reading the full cluster against the §3 mechanism stack and the §9
four-property criterion, four claims are now well-supported by the
convergent evidence:

\begin{enumerate}
\def\labelenumi{\arabic{enumi}.}
\item
  \textbf{The scaffold-versus-weights decomposition is independently
  confirmed.} GEPA, Meta-Harness, MOSS, RCL, and PACE all begin from the
  observation that model weights and the external artifact layer are
  distinct adaptation surfaces with different cost, reversibility, and
  scope properties. No surveyed paper disputes this decomposition; the
  variation is in which substrate each paper targets (skill document,
  harness code, system prompt, shared memory).
\item
  \textbf{Credit assignment is the central engineering bottleneck.} AHE,
  CANTANTE, Meta-Harness, and \citet{ong2026harnesseval} all identify
  credit assignment --- attributing a system-level failure to a specific
  scaffold coordinate --- as the problem that distinguishes disciplined
  Loop-1 evolution from trial-and-error. This is §3's operation 3, and
  it remains unsolved in the general multi-substrate case.
\item
  \textbf{Lifecycle hygiene is a precondition for descent, not an
  alternative to it.} Ratchet establishes that the gradient-analogue
  mechanism of §3 requires a maintained scaffold as a precondition. A
  dirty artifact pool (stale, bloated, miscurated) prevents the gate
  \(G_t\) from receiving a meaningful signal. Hygiene and optimisation
  are sequential dependencies, not competing paradigms.
\item
  \textbf{No surveyed system combines Loop-1 automation with
  cross-organisational Loop-2 governance.} SkillOpt has the strongest
  Loop-1 closure and the most systematic evidence-radius measurement,
  but it does not implement cross-organisational promotion, versioned
  lineage with RBAC, or rollback. CORAL implements asynchronous
  cross-agent promotion within one system but not across organisations.
  PACE implements two-timescale consolidation within one deployment.
  Meta-Harness implements Loop-1 closure over the harness substrate but
  no Loop-2 promotion mechanism. The four-property gap named in §9 is
  not reduced by any member of this cluster; it is, if anything,
  confirmed more strongly by the independent convergence of multiple
  research threads on the same single-organisation ceiling.
\end{enumerate}

The §3 picture is not aspirational. It is the limit that the field is
converging on from multiple directions simultaneously, under multiple
framings and for multiple substrates. What remains is not the
architecture --- that exists, in fragments, across SkillOpt, GEPA,
Meta-Harness, AHE, MOSS, CANTANTE, Ratchet, PACE, CORAL, and RCL --- but
the governance layer that would allow a locally-validated edit to be
promoted to a cross-organisational scope without losing the
auditability, rollback safety, and evidence-radius discipline that make
it trustworthy.

\subsubsection{G.10 SkillOpt against §9's four
properties}\label{g.10-skillopt-against-9s-four-properties}

Mapped onto §9's four properties of governed cross-tenant scaffold
optimization:

\begin{itemize}
\tightlist
\item
  \begin{enumerate}
  \def\labelenumi{(\alph{enumi})}
  \tightlist
  \item
    \textbf{Auto-gated Loop 1.} \emph{Satisfied.} The held-out
    selection-split gate at strictly-greater margin instantiates
    \(G_t\).
  \end{enumerate}
\item
  \begin{enumerate}
  \def\labelenumi{(\alph{enumi})}
  \setcounter{enumi}{1}
  \tightlist
  \item
    \textbf{Multi-tenant cross-org promotion.} \emph{Absent.} SkillOpt
    is a single-organisation research system; there is no governed
    cross-tenant aggregator.
  \end{enumerate}
\item
  \begin{enumerate}
  \def\labelenumi{(\alph{enumi})}
  \setcounter{enumi}{2}
  \tightlist
  \item
    \textbf{Versioned lineage with RBAC.} \emph{Partial.} The exported
    \texttt{best\_skill.md} is a versioned artifact, and the protected
    slow-update field separates fast and slow consolidation. There is no
    documented RBAC layer or audit-trail policy, and
    \texttt{edit\_apply\_report.json} is an operational log rather than
    an access-controlled lineage system.
  \end{enumerate}
\item
  \begin{enumerate}
  \def\labelenumi{(\alph{enumi})}
  \setcounter{enumi}{3}
  \tightlist
  \item
    \textbf{Rollback.} \emph{Absent.} The training loop selects the best
    validation-gated skill and exports it; no rollback envelope on
    deployed artifacts is reported.
  \end{enumerate}
\end{itemize}

The cross-model, cross-harness, and cross-benchmark transfer experiments
reported by SkillOpt are \emph{empirical evidence-radius measurements}:
they tell us that a skill optimised in one (model, benchmark, harness)
cell continues to improve performance when redeployed in nearby cells
without further optimisation. In §3 vocabulary, this is a measured
Loop-2 \emph{radius} for the artifact, not a \emph{governance mechanism}
for promoting the artifact. Both readings are useful: SkillOpt is the
strongest current Loop-1 proof and supplies the first systematic
transfer evidence at the scaffold layer, while the §9 cross-tenant
governance triad ((b), (c), (d)) remains unfilled.

\subsubsection{G.11 Cost and bloat}\label{g.11-cost-and-bloat}

SkillOpt reports 0.6M to 46.4M training tokens per absolute test-point
gain (their Table 6), with the procedural benchmarks (SpreadsheetBench,
OfficeQA, LiveMath) at the cheap end and search/visual benchmarks
(SearchQA, DocVQA) an order of magnitude more expensive per point. The
deployed artifact remains 379--1,995 tokens regardless. This is, to our
knowledge, the first published cost-per-point figure for scaffold-layer
descent, and it complements the §10 scaffold-bloat discussion (and the
Ratchet hygiene-precondition finding walked in G.7): the textual
learning rate \(L_t\) with cosine decay keeps the deployed skill compact
even when the optimiser explores many more edits than are accepted (only
1--4 edits survive into the final skill across the six benchmarks). The
mechanism that prevents the deployed artifact from bloating is exactly
the complexity penalty \(\Omega(z)\) of §3, instantiated as a hard cap
on accepted edits per step rather than a soft regulariser.

\begin{center}\rule{0.5\linewidth}{0.5pt}\end{center}


\nocite{*}
\bibliographystyle{plainnat}
\bibliography{references}

% NeurIPS paper checklist suppressed in this format-ready artifact.
% Uncomment the two lines below at actual submission time (the
% checklist file `checklist.tex` is filled and kept in the repo).
% \newpage
% \input{checklist.tex}

\end{document}
